\ifdefined\UOMTransportFragmentLoaded\else
\documentclass[11pt]{article}
\usepackage[margin=1in]{geometry}
\usepackage{amsmath,amssymb,amsthm,mathtools}
\usepackage{bm}
\usepackage{microtype}
\usepackage{enumitem}
\usepackage{booktabs}
\usepackage{array}
\usepackage{xcolor}
\usepackage{xr-hyper}
\usepackage{hyperref}
\usepackage[nameinlink,noabbrev]{cleveref}

\providecommand{\Fsloc}{\mathcal F_{\rm s.loc}}
\providecommand{\Qt}{Q_t}
\externaldocument[main-]{UOM/main}
\externaldocument[uomdesc-]{UOM/uom_descendant_note}
\externaldocument{UOM/completion_sector_split_model_note}

\hypersetup{
  colorlinks=true,
  linkcolor=blue!60!black,
  citecolor=blue!60!black,
  urlcolor=blue!60!black
}

\theoremstyle{definition}
\newtheorem{definition}{Definition}[section]
\newtheorem{assumption}{Assumption}[section]
\theoremstyle{plain}
\newtheorem{theorem}[definition]{Theorem}
\newtheorem{lemma}[definition]{Lemma}
\newtheorem{proposition}[definition]{Proposition}
\newtheorem{corollary}[definition]{Corollary}
\theoremstyle{remark}
\newtheorem{remark}[definition]{Remark}
\crefname{definition}{definition}{definitions}
\Crefname{definition}{Definition}{Definitions}
\crefname{assumption}{assumption}{assumptions}
\Crefname{assumption}{Assumption}{Assumptions}
\crefname{theorem}{theorem}{theorems}
\Crefname{theorem}{Theorem}{Theorems}
\crefname{lemma}{lemma}{lemmas}
\Crefname{lemma}{Lemma}{Lemmas}
\crefname{proposition}{proposition}{propositions}
\Crefname{proposition}{Proposition}{Propositions}
\crefname{corollary}{corollary}{corollaries}
\Crefname{corollary}{Corollary}{Corollaries}
\crefname{remark}{remark}{remarks}
\Crefname{remark}{Remark}{Remarks}

\newcommand{\cH}{\mathcal H}
\newcommand{\cA}{\mathcal A}
\newcommand{\cD}{\mathcal D}
\newcommand{\cB}{\mathcal B}
\newcommand{\cC}{\mathcal C}
\newcommand{\cL}{\mathcal L}
\newcommand{\cK}{\mathcal K}
\newcommand{\cR}{\mathcal R}
\newcommand{\cG}{\mathcal G}
\newcommand{\Ztwo}{\mathbb Z_2}
\newcommand{\Sys}{\mathrm{Sys}}
\newcommand{\Time}{\mathrm{Time}}
\newcommand{\RG}{\mathrm{RG}}
\newcommand{\Tr}{\mathrm{Tr}}
\newcommand{\diag}{\mathrm{diag}}
\newcommand{\ad}{\mathrm{ad}}
\newcommand{\id}{\mathrm{id}}
\newcommand{\norm}[1]{\left\lVert #1\right\rVert}
\newcommand{\abs}[1]{\left\lvert #1\right\rvert}
\newcommand{\ip}[2]{\left\langle #1,#2\right\rangle}

\externaldocument{UOM/compressed_log_spectral_selection_note}
\externaldocument{UOM/odd_selector_note}
\externaldocument{UOM/canonical_odd_interface_note}
\externaldocument[splitlift-]{UOM/split_controlled_completion_lift_note}
\title{Finite-Cutoff Active Response Transport Theorem}
\author{}
\date{}
\begin{document}
\maketitle
\fi

\section{Finite-Cutoff Active Response Transport Theorem}
\label{sec:active-response-transport-program}

This section assembles a finite-cutoff transport theorem for the odd descendant / Hamiltonization
lane of the 2T Main assumption package Definition~\ref{main-def:uom-main-theorem-assumptions}.
Accordingly, the phrase ``accepted tick'' remains literal UOM cutoff language throughout this note:
the transport theorem is pointwise on the accepted branch at fixed \(\Lambda_k\).

\begin{remark}[Physical odd sector and branchwise ambient-space conventions]
\label{rem:transport-conventions}
Throughout this note, the superscript \({}^-\) refers to physical \(\Gamma\)-odd parity on the
characterization layer, not to cohomological oddness.
Every ambiguity direction appearing below is therefore understood after passage to the physically
odd characterization sector \(E^-_{\Lambda_k}\).

All branchwise continuity statements are made after the standard accepted-branch identifications of
the relevant finite-cutoff source spaces, covariance-tangent spaces, and physical subspaces with
fixed ambient finite-cutoff spaces.
Operator-norm and Fisher-norm continuity are always understood with respect to those identifications.
\end{remark}

\begin{definition}[Odd ambiguity subspace on the characterization layer]
\label{def:transport-ambiguity-subspace}
For each accepted tick \(k\), let
\[
\mathcal A^-_{\Lambda_k}\subset E^-_{\Lambda_k}
\]
denote the closed linear span of the odd characterization-layer projections of the descendant
ambiguity directions
\[
\big(Q_tX+Q_uY\big)^-,\qquad \big(\partial_\mu K^\mu\big)^-,\qquad \Delta_{\rm ct}^{-},
\]
where \(K^\mu\) is boundary-admissible on the characterization layer and
\(\Delta_{\rm ct}^{-}\) is admissible in the sense of \cref{def:ambiguity-ct-char}.
Equivalently, \(\mathcal A^-_{\Lambda_k}\) is the odd ambiguity space already placed inside
\(\mathcal T^{\mathrm{op}}_{\Lambda_k}\) by \cref{cor:ambiguity-char}.
\end{definition}

\begin{definition}[Dominant compressed response line and scalar coefficient]
\label{def:dominant-response-functional}
Fix an accepted tick \(k\). Let
\[
u^{\mathrm{LoG}}_{\Lambda_k}
\]
and
\[
P_{-,\Lambda_k}^{\mathrm{LoG}}
\]
be the canonically selected compressed odd line and projector furnished by
\cref{cor:compressed-export-selection}. Define the physical covariance tangent induced by that line by
\[
\delta C^{\mathrm{LoG}}_{\Lambda_k}
:=
\Pi_{\mathrm{phys},\Lambda_k}\,\mathsf T_{\Lambda_k}[u^{\mathrm{LoG}}_{\Lambda_k}],
\]
and, whenever this tangent is nonzero, write
\[
L^{\mathrm{dom}}_{\Lambda_k}
:=
\mathbb R\,\widehat{\delta C}^{\mathrm{LoG}}_{\Lambda_k}
\]
for its Fisher-normalized line.
The raw dominant-response coefficient is
\begin{equation}
\widetilde{\rho}^{\mathrm{dom}}_{\Lambda_k}(F)
:=
\big\langle\!\big\langle
\Pi_{\mathrm{phys},\Lambda_k}\mathcal R_{\Lambda_k}(F),\,
\delta C^{\mathrm{LoG}}_{\Lambda_k}
\big\rangle\!\big\rangle_{C_{\ast,\Lambda_k}},
\qquad
F\in E^-_{\Lambda_k}.
\label{eq:rho-dom-raw-def}
\end{equation}
Whenever \(\delta C^{\mathrm{LoG}}_{\Lambda_k}\neq 0\), equivalently whenever dominant-line
realization holds at \(\Lambda_k\), define the normalized dominant-response coefficient by
\begin{equation}
\rho^{\mathrm{dom}}_{\Lambda_k}(F)
:=
\kappa_{\Lambda_k}^{-1/2}\,
\widetilde{\rho}^{\mathrm{dom}}_{\Lambda_k}(F),
\qquad
F\in E^-_{\Lambda_k}.
\label{eq:rho-dom-def}
\end{equation}
\end{definition}

\begin{proposition}[Well-posedness and continuity of the dominant-response coefficients]
\label{prop:rho-dom-continuity}
Fix an accepted tick \(k\).
Then the raw dominant-response coefficient
\[
\widetilde{\rho}^{\mathrm{dom}}_{\Lambda_k}:E^-_{\Lambda_k}\to\mathbb R
\]
is continuous and linear.
If dominant-line realization holds at \(\Lambda_k\), then the normalized dominant-response
coefficient
\[
\rho^{\mathrm{dom}}_{\Lambda_k}:E^-_{\Lambda_k}\to\mathbb R
\]
is also continuous and linear.
\end{proposition}

\begin{proof}
By \cref{prop:intrinsic-cov-response}, the intrinsic response map
\(
\mathcal R_{\Lambda_k}:E^-_{\Lambda_k}\to\mathcal T^{\mathrm{cov}}_{\Lambda_k}
\)
is continuous and linear.
Pairing its physical projection against the fixed covariance tangent
\(
\delta C^{\mathrm{LoG}}_{\Lambda_k}
\)
in the Fisher metric therefore gives a continuous linear functional, which is exactly
\(\widetilde{\rho}^{\mathrm{dom}}_{\Lambda_k}\).
If dominant-line realization holds, then \(\kappa_{\Lambda_k}>0\) by
\cref{prop:dominant-line-equivalent}, so \(\rho^{\mathrm{dom}}_{\Lambda_k}\) is obtained from
\(\widetilde{\rho}^{\mathrm{dom}}_{\Lambda_k}\) by multiplication by the fixed scalar
\(\kappa_{\Lambda_k}^{-1/2}\).
\end{proof}

\subsection*{Dominant-line realization}
\label{subsec:dominant-line-realization}

The compressed spectral-selection note already supplies the one-dimensional source line
\(\mathcal H^{\mathrm{LoG}}_{\Lambda_k}\).
The present transport note needs only the nonvanishing and branchwise continuity of its
\emph{physical} covariance push-forward.

\begin{definition}[Compressed-line transversality]
\label{def:compressed-line-transversality}
For an accepted tick \(k\), say that the welded response is \emph{transversely physical on the
compressed line} if
\begin{equation}
\Pi_{\mathrm{phys},\Lambda_k}\,\mathsf T_{\Lambda_k}\!\restriction_{\mathcal H^{\mathrm{LoG}}_{\Lambda_k}}
\neq 0.
\label{eq:compressed-transversality}
\end{equation}
Equivalently, the unique generator \(u^{\mathrm{LoG}}_{\Lambda_k}\) of the compressed odd line has
nonzero physical covariance response.
\end{definition}

\begin{definition}[Dominant-response norm]
\label{def:dominant-response-norm}
Whenever \(\delta C^{\mathrm{LoG}}_{\Lambda_k}\) is defined, set
\begin{equation}
\kappa_{\Lambda_k}
:=
\big\langle\!\big\langle
\delta C^{\mathrm{LoG}}_{\Lambda_k},\,
\delta C^{\mathrm{LoG}}_{\Lambda_k}
\big\rangle\!\big\rangle_{C_{\ast,\Lambda_k}}.
\label{eq:kappa-dom-def}
\end{equation}
This is the Fisher squared norm of the physical covariance tangent induced by the compressed welded
LoG source line.
\end{definition}

\begin{proposition}[Equivalent formulations of dominant-line realization]
\label{prop:dominant-line-equivalent}
For a fixed accepted tick \(k\), the following are equivalent.
\begin{enumerate}[leftmargin=2em]
\item the compressed welded LoG line is transversely physical in the sense of
  \cref{def:compressed-line-transversality};
\item the distinguished physical covariance tangent is nonzero,
  \[
  \delta C^{\mathrm{LoG}}_{\Lambda_k}\neq 0;
  \]
\item the dominant-response norm is strictly positive,
  \[
  \kappa_{\Lambda_k}>0;
  \]
\item there exists a physical covariance tangent
  \(
  B_{\Lambda_k}\in \mathcal T^{\mathrm{phys}}_{\Lambda_k}
  \)
  such that
  \[
  \big\langle\!\big\langle
  \delta C^{\mathrm{LoG}}_{\Lambda_k},\,B_{\Lambda_k}
  \big\rangle\!\big\rangle_{C_{\ast,\Lambda_k}}
  \neq 0;
  \]
\item in any Gaussian completion realization of
  \cref{cor:compressed-export-selection-completion}, the completion-sector covariance tangent
  representing \(\delta C^{\mathrm{LoG}}_{\Lambda_k}\) is nonzero.
\end{enumerate}
\end{proposition}

\begin{proof}
The equivalence of (1) and (2) is immediate from
\eqref{eq:compressed-transversality} and the definition
\[
\delta C^{\mathrm{LoG}}_{\Lambda_k}
=
\Pi_{\mathrm{phys},\Lambda_k}\mathsf T_{\Lambda_k}[u^{\mathrm{LoG}}_{\Lambda_k}],
\]
because \(\mathcal H^{\mathrm{LoG}}_{\Lambda_k}\) is the one-dimensional line generated by
\(u^{\mathrm{LoG}}_{\Lambda_k}\).

The equivalence of (2) and (3) follows from positive definiteness of the Fisher metric on the
covariance tangent space: a vector has positive Fisher norm if and only if it is nonzero.

Statement (3) implies (4) by taking
\(
B_{\Lambda_k}:=\delta C^{\mathrm{LoG}}_{\Lambda_k}
\),
which belongs to the physical tangent space by construction.
Conversely, if (4) holds then \(\delta C^{\mathrm{LoG}}_{\Lambda_k}\) cannot vanish, so (2) holds.

Finally, \cref{prop:intrinsic-cov-response,cor:compressed-export-selection-completion} identify the
abstract covariance tangent with any Gaussian completion representative of the same tangent object.
Therefore the abstract tangent is nonzero exactly when its completion realization is nonzero, giving
the equivalence of (2) and (5).
\end{proof}

\begin{theorem}[Branchwise realization and persistence of the dominant line]
\label{thm:dominant-line-persistence}
Let \(I\) be a connected accepted-branch segment in cutoff space.
Assume that:
\begin{enumerate}[leftmargin=2em]
\item the compressed projector \(P_{-,\Lambda}^{\mathrm{LoG}}\) varies continuously on \(I\);
\item the physical welded response operator
  \[
  \Lambda\longmapsto \Pi_{\mathrm{phys},\Lambda}\mathsf T_{\Lambda}
  \]
  is continuous on \(I\) in the operator topology induced by the Fisher norm on covariance
  tangents;
\item there exists a constant \(c_I>0\) such that
  \[
  \kappa_{\Lambda}\ge c_I
  \qquad
  \forall \Lambda\in I.
  \]
\end{enumerate}
Then:
\begin{enumerate}[leftmargin=2em]
\item \(\delta C^{\mathrm{LoG}}_{\Lambda}\neq 0\) for every \(\Lambda\in I\);
\item the normalized dominant vector
  \[
  \widehat{\delta C}^{\mathrm{LoG}}_{\Lambda}
  :=
  \kappa_{\Lambda}^{-1/2}\,\delta C^{\mathrm{LoG}}_{\Lambda}
  \]
  is continuous on \(I\);
\item the dominant physical line
  \[
  L^{\mathrm{dom}}_{\Lambda}
  =
  \mathbb R\,\widehat{\delta C}^{\mathrm{LoG}}_{\Lambda}
  \]
  is canonically transported along \(I\);
\item the Fisher-orthogonal rank-one projector
  \[
  \Pi^{\mathrm{dom}}_{\Lambda}
  :=
  \widehat{\delta C}^{\mathrm{LoG}}_{\Lambda}
  \otimes
  \widehat{\delta C}^{\mathrm{LoG}}_{\Lambda}
  \]
  depends continuously on \(\Lambda\in I\).
\end{enumerate}
\end{theorem}

\begin{proof}
By \cref{cor:compressed-export-selection}, the compressed projector
\(
P_{-,\Lambda}^{\mathrm{LoG}}
\)
and hence the normalized source vector \(u^{\mathrm{LoG}}_{\Lambda}\) vary continuously on the branch.
By the assumed continuity of \(\Pi_{\mathrm{phys},\Lambda}\mathsf T_{\Lambda}\), the map
\[
\Lambda\longmapsto
\delta C^{\mathrm{LoG}}_{\Lambda}
=
\Pi_{\mathrm{phys},\Lambda}\mathsf T_{\Lambda}[u^{\mathrm{LoG}}_{\Lambda}]
\]
is continuous in the Fisher norm.

The lower bound \(\kappa_{\Lambda}\ge c_I>0\) implies
\(
\delta C^{\mathrm{LoG}}_{\Lambda}\neq 0
\)
for every \(\Lambda\in I\), proving item (1).
Since the normalization map \(v\mapsto v/\|v\|_{C_\ast}\) is continuous on the complement of the
zero section, continuity of \(\delta C^{\mathrm{LoG}}_{\Lambda}\) and the positive lower bound on its
norm imply continuity of \(\widehat{\delta C}^{\mathrm{LoG}}_{\Lambda}\), proving item (2).
Items (3) and (4) are immediate consequences of item (2).
\end{proof}

\begin{corollary}[Local persistence from one-tick transversality]
\label{cor:dominant-line-local-persistence}
Assume the continuity hypotheses of \cref{thm:dominant-line-persistence}.
If the compressed line is transversely physical at one accepted tick \(\Lambda_{k_0}\), then there
exists a neighborhood \(I_{k_0}\) of \(\Lambda_{k_0}\) on the same accepted branch such that
\[
\delta C^{\mathrm{LoG}}_{\Lambda}\neq 0
\qquad
\forall \Lambda\in I_{k_0}.
\]
Hence the dominant physical line is locally well defined and continuously transported near
\(\Lambda_{k_0}\).
\end{corollary}

\begin{proof}
By \cref{prop:dominant-line-equivalent}, transversality at \(\Lambda_{k_0}\) is equivalent to
\(
\kappa_{\Lambda_{k_0}}>0
\).
Continuity of \(\kappa_{\Lambda}\) then gives a neighborhood \(I_{k_0}\) and a constant
\(c_{k_0}>0\) such that
\(
\kappa_{\Lambda}\ge c_{k_0}
\)
for all \(\Lambda\in I_{k_0}\).
Apply \cref{thm:dominant-line-persistence}.
\end{proof}

\begin{remark}[Practical witness criteria]
\label{rem:dominant-line-witnesses}
For applications, any one of the following suffices to discharge dominant-line realization at a fixed
accepted tick:
\begin{enumerate}[leftmargin=2em]
\item a direct lower bound \(\kappa_{\Lambda_k}>0\);
\item a physical witness tangent \(B_{\Lambda_k}\in\mathcal T^{\mathrm{phys}}_{\Lambda_k}\) with
  nonzero Fisher pairing against \(\delta C^{\mathrm{LoG}}_{\Lambda_k}\);
\item a Gaussian completion calculation showing that the completion covariance of the realized active
  odd mode changes nontrivially under the compressed welded LoG source.
\end{enumerate}
Each witness proves the same theorem-level fact by \cref{prop:dominant-line-equivalent}.
\end{remark}

\subsection*{Descent to the ambiguity quotient}
\label{subsec:ambiguity-quotient-descent}

Once the dominant line has been realized, the next task is purely structural:
show that the dominant-response coefficient is insensitive to ambiguity moves.
This is the exact point at which the transported descendant class becomes an affine class in the
sense required by Main rather than a collection of representatives.

\begin{definition}[Dominant-response quotient and descended descendant point]
\label{def:dominant-response-quotient}
For each accepted tick \(k\), define the dominant-response quotient by
\[
\mathfrak D^-_{\Lambda_k}
:=
E^-_{\Lambda_k}/\mathcal A^-_{\Lambda_k},
\]
with quotient map
\[
q_{\Lambda_k}:E^-_{\Lambda_k}\to \mathfrak D^-_{\Lambda_k}.
\]
If
\[
[\Xi^{(1)}_{\Lambda_k}]
=
J_{0,\Lambda_k}+\mathcal A^-_{\Lambda_k},
\]
then its image in the quotient is a single point,
\[
[\![\Xi^{(1)}_{\Lambda_k}]\!]
:=
q_{\Lambda_k}\!\big(J_{0,\Lambda_k}\big)
\in \mathfrak D^-_{\Lambda_k},
\]
independent of the chosen representative \(J_{0,\Lambda_k}\).
\end{definition}

\begin{proposition}[Universal descent criterion for the dominant-response coefficient]
\label{prop:dominant-quotient-universal}
Assume dominant-line realization at an accepted tick \(k\), so that
\(\rho^{\mathrm{dom}}_{\Lambda_k}\) is defined by \eqref{eq:rho-dom-def}.
Then the following are equivalent.
\begin{enumerate}[leftmargin=2em]
\item there exists a unique continuous linear functional
  \[
  \overline{\rho}^{\mathrm{dom}}_{\Lambda_k}:\mathfrak D^-_{\Lambda_k}\to\mathbb R
  \]
  such that
  \[
  \rho^{\mathrm{dom}}_{\Lambda_k}
  =
  \overline{\rho}^{\mathrm{dom}}_{\Lambda_k}\circ q_{\Lambda_k};
  \]
\item the ambiguity subspace is contained in the kernel:
  \[
  \mathcal A^-_{\Lambda_k}\subset \ker\rho^{\mathrm{dom}}_{\Lambda_k};
  \]
\item the dominant-response coefficient is constant on each ambiguity coset:
  \[
  \rho^{\mathrm{dom}}_{\Lambda_k}(F+A)
  =
  \rho^{\mathrm{dom}}_{\Lambda_k}(F)
  \qquad
  \forall F\in E^-_{\Lambda_k},\ \forall A\in\mathcal A^-_{\Lambda_k}.
  \]
\end{enumerate}
\end{proposition}

\begin{proof}
The equivalence of (2) and (3) is immediate from linearity:
\[
\rho^{\mathrm{dom}}_{\Lambda_k}(F+A)-\rho^{\mathrm{dom}}_{\Lambda_k}(F)
=
\rho^{\mathrm{dom}}_{\Lambda_k}(A).
\]

If (2) holds, define
\[
\overline{\rho}^{\mathrm{dom}}_{\Lambda_k}\big(q_{\Lambda_k}(F)\big)
:=
\rho^{\mathrm{dom}}_{\Lambda_k}(F).
\]
This is well defined precisely because changing representatives by an element of
\(\mathcal A^-_{\Lambda_k}\) does not change the value of \(\rho^{\mathrm{dom}}_{\Lambda_k}\).
Linearity and continuity are inherited from \(\rho^{\mathrm{dom}}_{\Lambda_k}\), and uniqueness
follows because \(q_{\Lambda_k}\) is surjective.

Conversely, if (1) holds and \(A\in\mathcal A^-_{\Lambda_k}\), then \(q_{\Lambda_k}(A)=0\), hence
\[
\rho^{\mathrm{dom}}_{\Lambda_k}(A)
=
\overline{\rho}^{\mathrm{dom}}_{\Lambda_k}\!\big(q_{\Lambda_k}(A)\big)
=
\overline{\rho}^{\mathrm{dom}}_{\Lambda_k}(0)
=
0.
\]
Thus (1) implies (2).
\end{proof}

\begin{corollary}[Constancy on the transported descendant affine class]
\label{cor:dominant-quotient-constancy}
Assume \cref{prop:dominant-quotient-universal}. Then
\[
\rho^{\mathrm{dom}}_{\Lambda_k}(J)
=
\overline{\rho}^{\mathrm{dom}}_{\Lambda_k}\big([\![\Xi^{(1)}_{\Lambda_k}]\!]\big)
\qquad
\forall J\in[\Xi^{(1)}_{\Lambda_k}].
\]
In particular, once the descended quotient functional exists, the descendant affine class carries a
single unambiguous dominant-response value.
\end{corollary}

\begin{proof}
Every \(J\in[\Xi^{(1)}_{\Lambda_k}]\) has the form
\(
J=J_{0,\Lambda_k}+A
\)
with \(A\in\mathcal A^-_{\Lambda_k}\), so
\[
q_{\Lambda_k}(J)=q_{\Lambda_k}(J_{0,\Lambda_k})=[\![\Xi^{(1)}_{\Lambda_k}]\!].
\]
Apply the factorization identity of \cref{prop:dominant-quotient-universal}.
\end{proof}

\begin{remark}[False stronger form and normalized special case]
\label{rem:dominant-quotient-false-stronger-form}
The quotient statement above is intrinsically a statement about the descended dominant scalar
\[
\overline{\rho}^{\mathrm{dom}}_{\Lambda_k}\big([\![\Xi^{(1)}_{\Lambda_k}]\!]\big),
\]
not about a bare \(\chi\)-compatibility condition.
At the weak transported-interface layer one has
\[
\chi_{\Lambda_k}
=
\nu_{\Lambda_k}^{-1}\rho^{\mathrm{dom}}_{\Lambda_k},
\]
so a representative-search condition of the form
\[
\rho^{\mathrm{dom}}_{\Lambda_k}(J)=s
\]
is equivalent to
\[
\chi_{\Lambda_k}(J)=\nu_{\Lambda_k}^{-1}s,
\]
not in general to \(\chi_{\Lambda_k}(J)=s\).
Only after an explicit normalization, for example \(\nu_{\Lambda_k}=1\), may the same scalar
comparison be rewritten in descended \(\bar\chi_{\Lambda_k}\) form without changing the value.
The finite-cutoff Wolfram audit contains exact countermodels to the stronger unscaled
\(\chi\)-level claim.
\end{remark}

\begin{proposition}[Generator-by-generator reduction]
\label{prop:dominant-quotient-generators}
Assume dominant-line realization at an accepted tick \(k\).
Suppose the dominant-response coefficient vanishes on the three generating ambiguity families:
\begin{enumerate}[leftmargin=2em]
\item for every odd-projected BRST/gauge-exact direction \(\big(Q_tX+Q_uY\big)^-\),
  \[
  \rho^{\mathrm{dom}}_{\Lambda_k}\!\big(\big(Q_tX+Q_uY\big)^-\big)=0;
  \]
\item for every odd-projected boundary-admissible divergence \(\big(\partial_\mu K^\mu\big)^-\),
  \[
  \rho^{\mathrm{dom}}_{\Lambda_k}\!\big(\big(\partial_\mu K^\mu\big)^-\big)=0;
  \]
\item for every admissible odd counterterm \(\Delta^-_{\rm ct}\),
  \[
  \rho^{\mathrm{dom}}_{\Lambda_k}(\Delta^-_{\rm ct})=0.
  \]
\end{enumerate}
Then
\[
\mathcal A^-_{\Lambda_k}\subset\ker\rho^{\mathrm{dom}}_{\Lambda_k},
\]
hence \(\rho^{\mathrm{dom}}_{\Lambda_k}\) descends uniquely to the quotient
\(\mathfrak D^-_{\Lambda_k}\).
\end{proposition}

\begin{proof}
By \cref{def:transport-ambiguity-subspace}, \(\mathcal A^-_{\Lambda_k}\) is the closed linear span of
the three families listed in the statement. Since \(\rho^{\mathrm{dom}}_{\Lambda_k}\) is linear, it
vanishes on the algebraic span of those generators. Since it is also continuous, it vanishes on the
closure of that span. Therefore every element of \(\mathcal A^-_{\Lambda_k}\) lies in the kernel.
Now apply \cref{prop:dominant-quotient-universal}.
\end{proof}

\begin{theorem}[Descent to the ambiguity quotient]
\label{thm:dominant-quotient-descent}
Fix an accepted tick \(k\) and assume dominant-line realization.
If the three vanishing statements of \cref{prop:dominant-quotient-generators} hold, then there
exists a unique descended dominant
functional
\[
\overline{\rho}^{\mathrm{dom}}_{\Lambda_k}:\mathfrak D^-_{\Lambda_k}\to\mathbb R
\]
such that
\[
\rho^{\mathrm{dom}}_{\Lambda_k}
=
\overline{\rho}^{\mathrm{dom}}_{\Lambda_k}\circ q_{\Lambda_k},
\]
and the transported descendant affine class has a single well-defined quotient value
\[
\overline{\rho}^{\mathrm{dom}}_{\Lambda_k}\big([\![\Xi^{(1)}_{\Lambda_k}]\!]\big).
\]
\end{theorem}

\begin{proof}
Apply \cref{prop:dominant-quotient-generators,cor:dominant-quotient-constancy}.
\end{proof}

\subsection*{Nonzero overlap on the descendant class}
\label{subsec:descendant-overlap}

After dominant-line realization and quotient descent, the descendant class defines a single point in
the ambiguity quotient, and the dominant-response functional descends to that quotient. This
subsection studies whether that descended descendant point is \emph{active}, i.e.\ not annihilated
by the descended dominant functional.

\begin{definition}[Dominant covariance projector]
\label{def:dominant-cov-projector}
Assume dominant-line realization at an accepted tick \(k\).
Define the Fisher-orthogonal rank-one projector onto the dominant physical line by
\begin{equation}
\Pi^{\mathrm{dom}}_{\Lambda_k}(X)
:=
\big\langle\!\big\langle
X,\,
\widehat{\delta C}^{\mathrm{LoG}}_{\Lambda_k}
\big\rangle\!\big\rangle_{C_{\ast,\Lambda_k}}\,
\widehat{\delta C}^{\mathrm{LoG}}_{\Lambda_k},
\qquad
X\in\mathcal T^{\mathrm{cov}}_{\Lambda_k}.
\label{eq:Pi-dom-def}
\end{equation}
\end{definition}

\begin{proposition}[Representative-level witness criteria]
\label{prop:descendant-overlap-witness}
Assume dominant-line realization at an accepted tick \(k\).
For any odd generator \(F\in E^-_{\Lambda_k}\), the following are equivalent.
\begin{enumerate}[leftmargin=2em]
\item the dominant-response coefficient is nonzero:
  \[
  \rho^{\mathrm{dom}}_{\Lambda_k}(F)\neq 0;
  \]
\item the dominant projection of the physical covariance response is nonzero:
  \[
  \Pi^{\mathrm{dom}}_{\Lambda_k}\Pi_{\mathrm{phys},\Lambda_k}\mathcal R_{\Lambda_k}(F)\neq 0;
  \]
\item the physical covariance response has nonzero Fisher pairing with the dominant tangent:
  \[
  \big\langle\!\big\langle
  \Pi_{\mathrm{phys},\Lambda_k}\mathcal R_{\Lambda_k}(F),\,
  \delta C^{\mathrm{LoG}}_{\Lambda_k}
  \big\rangle\!\big\rangle_{C_{\ast,\Lambda_k}}
  \neq 0;
  \]
\item in any Gaussian completion realization of
  \cref{cor:compressed-export-selection-completion}, the completion representative of
  \(\Pi_{\mathrm{phys},\Lambda_k}\mathcal R_{\Lambda_k}(F)\) has nonzero component on the realized
  active odd line.
\end{enumerate}
\end{proposition}

\begin{proof}
By \cref{def:dominant-cov-projector},
\[
\Pi^{\mathrm{dom}}_{\Lambda_k}\Pi_{\mathrm{phys},\Lambda_k}\mathcal R_{\Lambda_k}(F)
=
\rho^{\mathrm{dom}}_{\Lambda_k}(F)\,
\widehat{\delta C}^{\mathrm{LoG}}_{\Lambda_k}.
\]
Since \(\widehat{\delta C}^{\mathrm{LoG}}_{\Lambda_k}\neq 0\), this proves the equivalence of (1)
and (2).

Statement (3) is equivalent to (1) because
\[
\big\langle\!\big\langle
\Pi_{\mathrm{phys},\Lambda_k}\mathcal R_{\Lambda_k}(F),\,
\delta C^{\mathrm{LoG}}_{\Lambda_k}
\big\rangle\!\big\rangle_{C_{\ast,\Lambda_k}}
=
\|\delta C^{\mathrm{LoG}}_{\Lambda_k}\|_{C_{\ast,\Lambda_k}}\,
\rho^{\mathrm{dom}}_{\Lambda_k}(F),
\]
and the prefactor is strictly positive by dominant-line realization.

Finally, by \cref{prop:intrinsic-cov-response,cor:compressed-export-selection-completion}, the
abstract covariance tangent and its Gaussian completion representative describe the same physical
response object in different coordinates. Therefore the abstract dominant projection is nonzero if
and only if the completion representative has nonzero component on the realized active line,
proving the equivalence of (2) and (4).
\end{proof}

\begin{definition}[Descendant overlap margin]
\label{def:descendant-overlap-margin}
Assume dominant-line realization and quotient descent at an accepted tick \(k\).
Define the descendant overlap margin by
\begin{equation}
\nu_{\Lambda_k}
:=
\left|
\overline{\rho}^{\mathrm{dom}}_{\Lambda_k}\big([\![\Xi^{(1)}_{\Lambda_k}]\!]\big)
\right|.
\label{eq:descendant-overlap-margin}
\end{equation}
We say that the transported descendant class is \emph{dominant-active} if
\(
\nu_{\Lambda_k}>0
\).
\end{definition}

\begin{proposition}[Equivalent formulations of nonzero overlap on the descendant class]
\label{prop:descendant-overlap-equivalent}
Assume dominant-line realization and quotient descent at an accepted tick \(k\).
Then the following are equivalent.
\begin{enumerate}[leftmargin=2em]
\item the transported descendant sector has nonzero overlap with the canonically selected compressed
  dominant odd line;
\item there exists one representative \(J^{\mathrm{desc}}_{\Lambda_k}\in[\Xi^{(1)}_{\Lambda_k}]\)
  such that
  \[
  \Pi^{\mathrm{dom}}_{\Lambda_k}\Pi_{\mathrm{phys},\Lambda_k}
  \mathcal R_{\Lambda_k}(J^{\mathrm{desc}}_{\Lambda_k})
  \neq 0;
  \]
\item every representative \(J\in[\Xi^{(1)}_{\Lambda_k}]\) has the same nonzero dominant-response
  value:
  \[
  \rho^{\mathrm{dom}}_{\Lambda_k}(J)
  =
  \overline{\rho}^{\mathrm{dom}}_{\Lambda_k}\big([\![\Xi^{(1)}_{\Lambda_k}]\!]\big)
  \neq 0;
  \]
\item the descended descendant point is not annihilated by the descended dominant functional:
  \[
  [\![\Xi^{(1)}_{\Lambda_k}]\!]\notin \ker\overline{\rho}^{\mathrm{dom}}_{\Lambda_k};
  \]
\item the descendant overlap margin is strictly positive:
  \[
  \nu_{\Lambda_k}>0.
  \]
\end{enumerate}
\end{proposition}

\begin{proof}
The equivalence of (1) and (2) follows immediately from
\cref{prop:descendant-overlap-witness}.

If (2) holds, then
\(
\rho^{\mathrm{dom}}_{\Lambda_k}(J^{\mathrm{desc}}_{\Lambda_k})\neq 0
\)
by \cref{prop:descendant-overlap-witness}. By
\cref{cor:dominant-quotient-constancy}, every other representative of
\([\Xi^{(1)}_{\Lambda_k}]\) has the same value, proving (3).
Conversely, (3) obviously implies (1).

The equivalence of (3) and (4) is immediate from the identity
\[
\rho^{\mathrm{dom}}_{\Lambda_k}(J)
=
\overline{\rho}^{\mathrm{dom}}_{\Lambda_k}\big([\![\Xi^{(1)}_{\Lambda_k}]\!]\big)
\qquad
\forall J\in[\Xi^{(1)}_{\Lambda_k}]
\]
of \cref{cor:dominant-quotient-constancy}.

Finally, (4) and (5) are equivalent by the definition of \(\nu_{\Lambda_k}\) in
\cref{def:descendant-overlap-margin}.
\end{proof}

\begin{theorem}[Nonzero overlap on the descendant class]
\label{thm:descendant-overlap}
Fix an accepted tick \(k\), and assume dominant-line realization together with quotient descent.
If any one of the equivalent conditions in \cref{prop:descendant-overlap-equivalent} holds, then the
following conclusions obtain.
In particular, it is enough to construct a single representative
\[
J^{\mathrm{desc}}_{\Lambda_k}\in[\Xi^{(1)}_{\Lambda_k}]
\]
whose dominant projected covariance response is nonzero.

When this holds, the quotient descendant point is dominant-active, the scalar value
\[
\overline{\rho}^{\mathrm{dom}}_{\Lambda_k}\big([\![\Xi^{(1)}_{\Lambda_k}]\!]\big)
\]
is nonzero, and the sign
\begin{equation}
\varepsilon_{\Lambda_k}
:=
\operatorname{sgn}\!\Big(
\overline{\rho}^{\mathrm{dom}}_{\Lambda_k}\big([\![\Xi^{(1)}_{\Lambda_k}]\!]\big)
\Big)
\in\{\pm1\}
\label{eq:epsilon-from-overlap}
\end{equation}
is well defined.
Thus the transported descendant sector has nonzero overlap with the canonically selected compressed
dominant odd line.
\end{theorem}

\begin{proof}
The equivalence package is exactly \cref{prop:descendant-overlap-equivalent}.
Once one representative has nonzero dominant projected response, the same proposition gives the
nonzero scalar value on the descended descendant point. The sign in \eqref{eq:epsilon-from-overlap}
is therefore well defined.
\end{proof}

\subsection*{Conditional mod-null reentry on the selected scalar line}
\label{subsec:conditional-mod-null-reentry}

\begin{definition}[Parity-stable mod-null subspace transverse to the selected scalar line]
\label{def:mod-null-transverse-datum}
Fix an accepted tick \(k\), assume dominant-line realization and quotient descent, and let
\[
\overline L^{\mathrm{Main}}_{\Lambda_k}
=
\mathbb R\,\overline\ell^{\mathrm{Main}}_{\Lambda_k}
\subset
\mathfrak D^-_{\Lambda_k}
\]
be a one-dimensional selected scalar line in the ambiguity quotient.
A \emph{parity-stable mod-null datum transverse to}
\(
\overline L^{\mathrm{Main}}_{\Lambda_k}
\)
consists of a linear subspace
\[
N^{\mathrm{mn}}_{\Lambda_k}\subset \mathfrak D^-_{\Lambda_k}
\]
such that:
\begin{enumerate}[leftmargin=2em,label=(\alph*)]
\item \(N^{\mathrm{mn}}_{\Lambda_k}\) is stable under the descended physical parity on
  \(\mathfrak D^-_{\Lambda_k}\);
\item the selected scalar line is transverse to \(N^{\mathrm{mn}}_{\Lambda_k}\):
  \[
  \overline L^{\mathrm{Main}}_{\Lambda_k}\cap N^{\mathrm{mn}}_{\Lambda_k}=\{0\}.
  \]
\end{enumerate}
Define
\[
\mathcal V^{1,\mathrm{mn}}_{\Lambda_k}
:=
\mathfrak D^-_{\Lambda_k}/N^{\mathrm{mn}}_{\Lambda_k}
\]
with quotient map
\[
\pi^{\mathrm{mn}}_{\Lambda_k}:
\mathfrak D^-_{\Lambda_k}\to \mathcal V^{1,\mathrm{mn}}_{\Lambda_k}.
\]
Because \(\mathfrak D^-_{\Lambda_k}\) is itself a quotient of the physical \(\Gamma\)-odd sector,
the descended parity acts on \(\mathfrak D^-_{\Lambda_k}\) by \(-\mathrm{id}\).
Thus the parity-stability clause yields a descended involution
\[
\Gamma^{\mathrm{mn}}_{\Lambda_k}
\bigl(
\pi^{\mathrm{mn}}_{\Lambda_k}(x)
\bigr)
:=
\pi^{\mathrm{mn}}_{\Lambda_k}(-x),
\qquad
x\in \mathfrak D^-_{\Lambda_k}.
\]
\end{definition}

\begin{definition}[Dominant-null parity-stable mod-null datum]
\label{def:dominant-null-mod-null-datum}
Fix an accepted tick \(k\), assume dominant-line realization and quotient descent, and let
\[
\overline L^{\mathrm{Main}}_{\Lambda_k}
=
\mathbb R\,\overline\ell^{\mathrm{Main}}_{\Lambda_k}
\subset
\mathfrak D^-_{\Lambda_k}
\]
be a one-dimensional selected scalar line in the ambiguity quotient.
A parity-stable mod-null datum
\[
N^{\mathrm{mn}}_{\Lambda_k}\subset \mathfrak D^-_{\Lambda_k}
\]
is called \emph{dominant-null} if it is contained in the kernel of the descended dominant
functional:
\[
N^{\mathrm{mn}}_{\Lambda_k}\subset \ker \overline{\rho}^{\mathrm{dom}}_{\Lambda_k}.
\]
It is understood that
\(
N^{\mathrm{mn}}_{\Lambda_k}
\)
is stable under the descended physical parity on
\(
\mathfrak D^-_{\Lambda_k}
\),
and that the associated quotient
\(
\mathfrak D^-_{\Lambda_k}/N^{\mathrm{mn}}_{\Lambda_k}
\)
carries the descended involution from \cref{def:mod-null-transverse-datum}.
The canonical dominant-kernel choice is isolated separately in
\cref{def:canonical-dominant-null-datum}.
The transversality lemma below uses only this kernel-subordination property.
\end{definition}

\begin{definition}[Canonical dominant-null datum]
\label{def:canonical-dominant-null-datum}
Fix an accepted tick \(k\) and assume dominant-line realization together with quotient descent.
Define the \emph{canonical dominant-null datum} by
\[
N^{\mathrm{dom-null}}_{\Lambda_k}
:=
\ker \overline{\rho}^{\mathrm{dom}}_{\Lambda_k}
\subset
\mathfrak D^-_{\Lambda_k}.
\]
This datum depends only on the descended dominant functional on the ambiguity quotient.
No additional reduced-welded or completion-side choice is built into the definition.
\end{definition}

\begin{lemma}[Dominant-kernel criterion for selected-line transversality]
\label{lem:dominant-kernel-transversality}
Fix an accepted tick \(k\), assume dominant-line realization together with quotient descent, and let
\[
\overline L^{\mathrm{Main}}_{\Lambda_k}
=
\mathbb R\,\overline\ell^{\mathrm{Main}}_{\Lambda_k}
\subset
\mathfrak D^-_{\Lambda_k}
\]
be a one-dimensional selected scalar line with
\[
\overline{\rho}^{\mathrm{dom}}_{\Lambda_k}
\!\bigl(
\overline\ell^{\mathrm{Main}}_{\Lambda_k}
\bigr)
\neq 0.
\]
If
\[
N^{\mathrm{mn}}_{\Lambda_k}\subset \mathfrak D^-_{\Lambda_k}
\]
is a dominant-null parity-stable mod-null datum in the sense of
\cref{def:dominant-null-mod-null-datum},
then
\[
\overline L^{\mathrm{Main}}_{\Lambda_k}\cap N^{\mathrm{mn}}_{\Lambda_k}=\{0\}.
\]
So every dominant-null parity-stable mod-null datum is automatically transverse to the selected
scalar line.
\end{lemma}

\begin{proof}
Let
\(
x\in \overline L^{\mathrm{Main}}_{\Lambda_k}\cap N^{\mathrm{mn}}_{\Lambda_k}
\).
Because
\(
\overline L^{\mathrm{Main}}_{\Lambda_k}
=
\mathbb R\,\overline\ell^{\mathrm{Main}}_{\Lambda_k}
\),
there exists \(a\in\mathbb R\) with
\(
x=a\,\overline\ell^{\mathrm{Main}}_{\Lambda_k}
\).
Since
\(
N^{\mathrm{mn}}_{\Lambda_k}\subset
\ker \overline{\rho}^{\mathrm{dom}}_{\Lambda_k}
\),
we have
\[
0
=
\overline{\rho}^{\mathrm{dom}}_{\Lambda_k}(x)
=
a\,
\overline{\rho}^{\mathrm{dom}}_{\Lambda_k}
\!\bigl(
\overline\ell^{\mathrm{Main}}_{\Lambda_k}
\bigr).
\]
The coefficient
\(
\overline{\rho}^{\mathrm{dom}}_{\Lambda_k}
(\overline\ell^{\mathrm{Main}}_{\Lambda_k})
\)
is nonzero by hypothesis, so \(a=0\).
Therefore \(x=0\), proving
\(
\overline L^{\mathrm{Main}}_{\Lambda_k}\cap N^{\mathrm{mn}}_{\Lambda_k}=\{0\}
\).
\end{proof}

\begin{theorem}[The dominant kernel is the canonical dominant-null datum]
\label{thm:canonical-dominant-null-datum}
Fix an accepted tick \(k\), assume dominant-line realization together with quotient descent, and let
\[
\overline L^{\mathrm{Main}}_{\Lambda_k}
=
\mathbb R\,\overline\ell^{\mathrm{Main}}_{\Lambda_k}
\subset
\mathfrak D^-_{\Lambda_k}
\]
be a one-dimensional selected scalar line with
\[
\overline{\rho}^{\mathrm{dom}}_{\Lambda_k}
\!\bigl(
\overline\ell^{\mathrm{Main}}_{\Lambda_k}
\bigr)
\neq 0.
\]
Define
\[
N^{\mathrm{dom-null}}_{\Lambda_k}
:=
\ker \overline{\rho}^{\mathrm{dom}}_{\Lambda_k}
\subset
\mathfrak D^-_{\Lambda_k}.
\]
Then:
\begin{enumerate}[leftmargin=2em]
\item \(N^{\mathrm{dom-null}}_{\Lambda_k}\) is stable under the descended physical parity on
  \(\mathfrak D^-_{\Lambda_k}\);
\item \(N^{\mathrm{dom-null}}_{\Lambda_k}\) is a dominant-null parity-stable mod-null datum in the
  sense of \cref{def:dominant-null-mod-null-datum};
\item the selected scalar line is automatically transverse to the canonical dominant-null datum:
  \[
  \overline L^{\mathrm{Main}}_{\Lambda_k}
  \cap
  N^{\mathrm{dom-null}}_{\Lambda_k}
  =
  \{0\}.
  \]
\end{enumerate}
Hence the canonical dominant kernel is admissible as the mod-null datum for the conditional
reentry theorem.
\end{theorem}

\begin{proof}
Let \(x\in N^{\mathrm{dom-null}}_{\Lambda_k}\).
Because the descended physical parity acts on \(\mathfrak D^-_{\Lambda_k}\) by \(-\mathrm{id}\),
we have
\[
\overline{\rho}^{\mathrm{dom}}_{\Lambda_k}(-x)
=
-\overline{\rho}^{\mathrm{dom}}_{\Lambda_k}(x)
=
0.
\]
So \(-x\in N^{\mathrm{dom-null}}_{\Lambda_k}\), proving parity stability.
Item \textnormal{(2)} is then exactly \cref{def:dominant-null-mod-null-datum} with
\(
N^{\mathrm{mn}}_{\Lambda_k}=N^{\mathrm{dom-null}}_{\Lambda_k}
\).
Item \textnormal{(3)} is \cref{lem:dominant-kernel-transversality} applied to this canonical
choice.
\end{proof}

\begin{lemma}[Transversality reduces mod-null reentry to quotient construction]
\label{lem:mod-null-transversality-reduction}
Fix an accepted tick \(k\), assume dominant-line realization together with quotient descent, and let
\[
\overline L^{\mathrm{Main}}_{\Lambda_k}
=
\mathbb R\,\overline\ell^{\mathrm{Main}}_{\Lambda_k}
\subset
\mathfrak D^-_{\Lambda_k}
\]
be a one-dimensional selected scalar line with
\[
\overline{\rho}^{\mathrm{dom}}_{\Lambda_k}
\!\bigl(
\overline\ell^{\mathrm{Main}}_{\Lambda_k}
\bigr)
\neq 0.
\]
Assume moreover a parity-stable mod-null datum transverse to
\(
\overline L^{\mathrm{Main}}_{\Lambda_k}
\)
in the sense of \cref{def:mod-null-transverse-datum}.
Define
\[
\ell^{\mathrm{sel,mn}}_{\Lambda_k}
:=
\pi^{\mathrm{mn}}_{\Lambda_k}
\!\bigl(
\overline\ell^{\mathrm{Main}}_{\Lambda_k}
\bigr)
\]
and
\[
L^{\mathrm{sel,mn}}_{\Lambda_k}
:=
\mathbb R\,\ell^{\mathrm{sel,mn}}_{\Lambda_k}
\subset
\mathcal V^{1,\mathrm{mn}}_{\Lambda_k}.
\]
Then:
\begin{enumerate}[leftmargin=2em]
\item the restriction of
  \(
  \pi^{\mathrm{mn}}_{\Lambda_k}
  \)
  to
  \(
  \overline L^{\mathrm{Main}}_{\Lambda_k}
  \)
  is injective;
\item \(L^{\mathrm{sel,mn}}_{\Lambda_k}\) is a well-defined one-dimensional line in the mod-null
  quotient, generated by the nonzero class
  \(
  \ell^{\mathrm{sel,mn}}_{\Lambda_k}
  \);
\item the mod-null selected realization map
  \[
  \mathfrak R^{\mathrm{mn}}_{\Lambda_k}:
  L^{\mathrm{sel,mn}}_{\Lambda_k}
  \longrightarrow
  \mathbb R\,u^{\mathrm{LoG}}_{\Lambda_k},
  \qquad
  \mathfrak R^{\mathrm{mn}}_{\Lambda_k}
  \!\bigl(
  a\,\ell^{\mathrm{sel,mn}}_{\Lambda_k}
  \bigr)
  :=
  a\,
  \overline{\rho}^{\mathrm{dom}}_{\Lambda_k}
  \!\bigl(
  \overline\ell^{\mathrm{Main}}_{\Lambda_k}
  \bigr)\,
  u^{\mathrm{LoG}}_{\Lambda_k}
  \]
  is well defined and faithful;
\item the descended involution acts by sign:
  \[
  \Gamma^{\mathrm{mn}}_{\Lambda_k}
  \!\bigl(
  \ell^{\mathrm{sel,mn}}_{\Lambda_k}
  \bigr)
  =
  -
  \ell^{\mathrm{sel,mn}}_{\Lambda_k}.
  \]
\end{enumerate}
In particular, once the transversality condition
\(
\overline L^{\mathrm{Main}}_{\Lambda_k}\cap N^{\mathrm{mn}}_{\Lambda_k}=\{0\}
\)
is known, the remaining mod-null reentry step reduces to quotient construction plus parity descent.
\end{lemma}

\begin{proof}
If
\(
x\in \overline L^{\mathrm{Main}}_{\Lambda_k}
\)
and
\(
\pi^{\mathrm{mn}}_{\Lambda_k}(x)=0
\),
then
\(
x\in N^{\mathrm{mn}}_{\Lambda_k}
\)
by definition of the quotient kernel.
Hence
\(
x\in
\overline L^{\mathrm{Main}}_{\Lambda_k}\cap N^{\mathrm{mn}}_{\Lambda_k}
=
\{0\}
\),
which proves injectivity.
Since
\(
\overline\ell^{\mathrm{Main}}_{\Lambda_k}\neq 0
\),
its image
\(
\ell^{\mathrm{sel,mn}}_{\Lambda_k}
\)
is therefore nonzero, and the image of a one-dimensional line under an injective linear map is
again one-dimensional.
This proves items \textnormal{(1)} and \textnormal{(2)}.

Injectivity on
\(
\overline L^{\mathrm{Main}}_{\Lambda_k}
\)
implies that every element of
\(
L^{\mathrm{sel,mn}}_{\Lambda_k}
\)
has a unique expression
\(
a\,\ell^{\mathrm{sel,mn}}_{\Lambda_k}
\),
so the displayed formula defines
\(
\mathfrak R^{\mathrm{mn}}_{\Lambda_k}
\)
unambiguously.
Its coefficient
\(
\overline{\rho}^{\mathrm{dom}}_{\Lambda_k}
(\overline\ell^{\mathrm{Main}}_{\Lambda_k})
\)
is nonzero by hypothesis, and
\(
u^{\mathrm{LoG}}_{\Lambda_k}\neq 0
\)
by \cref{cor:compressed-export-selection}.
Therefore
\(
\mathfrak R^{\mathrm{mn}}_{\Lambda_k}
(a\,\ell^{\mathrm{sel,mn}}_{\Lambda_k})
=
0
\)
implies
\(
a=0
\),
so the map is faithful.
This proves item \textnormal{(3)}.

By construction of the descended involution,
\[
\Gamma^{\mathrm{mn}}_{\Lambda_k}
\!\bigl(
\ell^{\mathrm{sel,mn}}_{\Lambda_k}
\bigr)
=
\Gamma^{\mathrm{mn}}_{\Lambda_k}
\!\Bigl(
\pi^{\mathrm{mn}}_{\Lambda_k}
\!\bigl(
\overline\ell^{\mathrm{Main}}_{\Lambda_k}
\bigr)
\Bigr)
=
\pi^{\mathrm{mn}}_{\Lambda_k}
\!\bigl(
-\overline\ell^{\mathrm{Main}}_{\Lambda_k}
\bigr)
=
-\ell^{\mathrm{sel,mn}}_{\Lambda_k},
\]
which proves item \textnormal{(4)}.
\end{proof}

\begin{theorem}[Conditional mod-null reentry on the selected scalar line]
\label{thm:conditional-mod-null-reentry}
Fix an accepted tick \(k\) and the next accepted receiver step \(k+1\).
Assume the hypotheses and notation of
\cref{splitlift-thm:splitlift-scalar-descendant-lift}.
In particular, let
\[
\overline L^{\mathrm{Main}}_{\Lambda_{k+1}}
=
\mathbb R\,\overline\ell^{\mathrm{Main}}_{\Lambda_{k+1}}
\subset
\mathfrak D^-_{\Lambda_{k+1}}
\]
be the nonzero selected scalar descendant line supplied there.
Assume moreover a parity-stable mod-null datum transverse to
\(
\overline L^{\mathrm{Main}}_{\Lambda_{k+1}}
\)
in the sense of \cref{def:mod-null-transverse-datum}.
Define
\[
\ell^{\mathrm{sel,mn}}_{\Lambda_{k+1}}
:=
\pi^{\mathrm{mn}}_{\Lambda_{k+1}}
\!\bigl(
\overline\ell^{\mathrm{Main}}_{\Lambda_{k+1}}
\bigr)
\]
and
\[
L^{\mathrm{sel,mn}}_{\Lambda_{k+1}}
:=
\mathbb R\,\ell^{\mathrm{sel,mn}}_{\Lambda_{k+1}}
\subset
\mathcal V^{1,\mathrm{mn}}_{\Lambda_{k+1}}.
\]
Then:
\begin{enumerate}[leftmargin=2em]
\item \(L^{\mathrm{sel,mn}}_{\Lambda_{k+1}}\) is a well-defined one-dimensional selected descendant
  line in the mod-null quotient;
\item the induced mod-null selected realization map
  \[
  \mathfrak R^{\mathrm{mn}}_{\Lambda_{k+1}}:
  L^{\mathrm{sel,mn}}_{\Lambda_{k+1}}
  \longrightarrow
  \mathbb R\,u^{\mathrm{LoG}}_{\Lambda_{k+1}},
  \qquad
  \mathfrak R^{\mathrm{mn}}_{\Lambda_{k+1}}
  \!\bigl(
  a\,\ell^{\mathrm{sel,mn}}_{\Lambda_{k+1}}
  \bigr)
  :=
  a\,
  \overline{\rho}^{\mathrm{dom}}_{\Lambda_{k+1}}
  \!\bigl(
  \overline\ell^{\mathrm{Main}}_{\Lambda_{k+1}}
  \bigr)\,
  u^{\mathrm{LoG}}_{\Lambda_{k+1}}
  \]
  is well defined and faithful;
\item the descended involution acts by sign:
  \[
  \Gamma^{\mathrm{mn}}_{\Lambda_{k+1}}
  \!\bigl(
  \ell^{\mathrm{sel,mn}}_{\Lambda_{k+1}}
  \bigr)
  =
  -
  \ell^{\mathrm{sel,mn}}_{\Lambda_{k+1}}.
  \]
\end{enumerate}
\end{theorem}

\begin{proof}
By \cref{splitlift-thm:splitlift-scalar-descendant-lift},
\[
\overline\ell^{\mathrm{Main}}_{\Lambda_{k+1}}\neq 0
\qquad\text{and}\qquad
\overline{\rho}^{\mathrm{dom}}_{\Lambda_{k+1}}
\!\bigl(
\overline\ell^{\mathrm{Main}}_{\Lambda_{k+1}}
\bigr)
\neq 0.
\]
So every hypothesis of
\cref{lem:mod-null-transversality-reduction}
is in force at \(\Lambda_{k+1}\).
Applying that lemma gives items \textnormal{(1)}--\textnormal{(3)} immediately.
\end{proof}

\begin{corollary}[Canonical dominant-kernel discharge of the mod-null reentry step]
\label{cor:dominant-kernel-mod-null-reentry}
Fix an accepted tick \(k\) and the next accepted receiver step \(k+1\).
Assume the hypotheses and notation of
\cref{splitlift-thm:splitlift-scalar-descendant-lift}.
Define the canonical dominant-null datum
\[
N^{\mathrm{dom-null}}_{\Lambda_{k+1}}
:=
\ker \overline{\rho}^{\mathrm{dom}}_{\Lambda_{k+1}}
\subset
\mathfrak D^-_{\Lambda_{k+1}}.
\]
Then \(N^{\mathrm{dom-null}}_{\Lambda_{k+1}}\) is admissible as the mod-null datum of
\cref{thm:conditional-mod-null-reentry}, the transversality hypothesis of that theorem is
automatic, and hence the conclusions of that theorem hold with this canonical datum.
\end{corollary}

\begin{proof}
By \cref{splitlift-thm:splitlift-scalar-descendant-lift},
\[
\overline{\rho}^{\mathrm{dom}}_{\Lambda_{k+1}}
\!\bigl(
\overline\ell^{\mathrm{Main}}_{\Lambda_{k+1}}
\bigr)
\neq 0.
\]
So \cref{thm:canonical-dominant-null-datum} applies at \(\Lambda_{k+1}\).
It shows that
\(
N^{\mathrm{dom-null}}_{\Lambda_{k+1}}
\)
is a dominant-null parity-stable mod-null datum and that
\[
\overline L^{\mathrm{Main}}_{\Lambda_{k+1}}
\cap
N^{\mathrm{dom-null}}_{\Lambda_{k+1}}
=
\{0\}.
\]
This is exactly the admissibility and transversality clause required by
\cref{thm:conditional-mod-null-reentry}, so that theorem now applies directly with the canonical
datum.
\end{proof}

\begin{remark}[Conditional character of the reentry step]
\label{rem:conditional-mod-null-reentry}
\cref{thm:conditional-mod-null-reentry} is a quotient-level reentry statement.
After
\cref{lem:dominant-kernel-transversality,thm:canonical-dominant-null-datum,cor:dominant-kernel-mod-null-reentry},
the canonical dominant-null datum is fixed as
\(
\ker\overline{\rho}^{\mathrm{dom}}_{\Lambda_{k+1}}
\)
on \(\mathfrak D^-_{\Lambda_{k+1}}\).
What remains is passage to the mod-null quotient together with parity descent on the odd line.
Accordingly, the theorem should not be read as a statement on the raw descendant ambiguity quotient
or on any stronger physical quotient.
\end{remark}

\subsection*{Branchwise orientation transport}
\label{subsec:branchwise-orientation}

This subsection studies transport of the descended descendant sign along a connected accepted
branch. The sign can be transported once the descended descendant value varies continuously with the
cutoff. The compressed spectral-selection theorem already gives continuity of the dominant
projector, and \cref{thm:dominant-line-persistence} gives continuity of the dominant physical line.
We combine those inputs with continuity of the projected descendant response.

\begin{definition}[Branchwise descendant overlap function]
\label{def:branchwise-overlap-function}
Let \(I\) be a connected accepted-branch segment.
Assume dominant-line realization and quotient descent for every \(\Lambda\in I\).
Suppose there is a branchwise choice of representatives
\[
J^{\mathrm{desc}}_{\Lambda}\in[\Xi^{(1)}_{\Lambda}]
\qquad
(\Lambda\in I)
\]
such that the projected covariance response
\[
\Lambda\longmapsto
\Pi_{\mathrm{phys},\Lambda}\mathcal R_{\Lambda}(J^{\mathrm{desc}}_{\Lambda})
\]
is continuous in the Fisher topology.
Define the associated branchwise descendant overlap function by
\begin{equation}
m_I(\Lambda)
:=
\rho^{\mathrm{dom}}_{\Lambda}(J^{\mathrm{desc}}_{\Lambda})
=
\overline{\rho}^{\mathrm{dom}}_{\Lambda}\big([\![\Xi^{(1)}_{\Lambda}]\!]\big).
\label{eq:branchwise-overlap-function}
\end{equation}
\end{definition}

\begin{proposition}[Continuity of the descended overlap scalar]
\label{prop:branchwise-overlap-continuity}
Let \(I\) be a connected accepted-branch segment.
Assume:
\begin{enumerate}[leftmargin=2em]
\item the hypotheses of \cref{thm:dominant-line-persistence} on \(I\);
\item quotient descent holds for every \(\Lambda\in I\);
\item the branchwise choice of representatives in
  \cref{def:branchwise-overlap-function} exists.
\end{enumerate}
Then the scalar function \(m_I:I\to\mathbb R\) is well defined, independent of the chosen
branchwise representatives, and continuous.
Moreover,
\begin{equation}
|m_I(\Lambda)|
=
\nu_{\Lambda}
\qquad
\forall \Lambda\in I.
\label{eq:branchwise-overlap-margin}
\end{equation}
\end{proposition}

\begin{proof}
Independence of the chosen branchwise representatives is exactly the coset constancy proved in
\cref{cor:dominant-quotient-constancy}: if
\(
J'^{\mathrm{desc}}_{\Lambda}-J^{\mathrm{desc}}_{\Lambda}\in\mathcal A^-_{\Lambda}
\),
then
\[
\rho^{\mathrm{dom}}_{\Lambda}(J'^{\mathrm{desc}}_{\Lambda})
=
\rho^{\mathrm{dom}}_{\Lambda}(J^{\mathrm{desc}}_{\Lambda})
=
\overline{\rho}^{\mathrm{dom}}_{\Lambda}\big([\![\Xi^{(1)}_{\Lambda}]\!]\big).
\]
Thus \(m_I\) is well defined.

By \cref{thm:dominant-line-persistence}, the normalized dominant vector
\(
\widehat{\delta C}^{\mathrm{LoG}}_{\Lambda}
\)
and the dominant projector
\(
\Pi^{\mathrm{dom}}_{\Lambda}
\)
depend continuously on \(\Lambda\in I\).
By assumption, the response vector
\[
X_{\Lambda}
:=
\Pi_{\mathrm{phys},\Lambda}\mathcal R_{\Lambda}(J^{\mathrm{desc}}_{\Lambda})
\]
is also continuous on \(I\).
Using \cref{def:dominant-cov-projector}, we have
\[
\Pi^{\mathrm{dom}}_{\Lambda}X_{\Lambda}
=
m_I(\Lambda)\,\widehat{\delta C}^{\mathrm{LoG}}_{\Lambda}.
\]
The left-hand side is continuous because both factors are continuous, while
\(
\widehat{\delta C}^{\mathrm{LoG}}_{\Lambda}\neq 0
\)
for all \(\Lambda\in I\) by \cref{thm:dominant-line-persistence}.
Therefore the scalar coefficient \(m_I(\Lambda)\) is continuous.

Finally, \eqref{eq:branchwise-overlap-margin} is just the definition of \(\nu_{\Lambda}\) from
\cref{def:descendant-overlap-margin} together with \eqref{eq:branchwise-overlap-function}.
\end{proof}

\begin{theorem}[Branchwise orientation transport on the descendant class]
\label{thm:branchwise-orientation-transport}
Let \(I\) be a connected accepted-branch segment and assume the hypotheses of
\cref{prop:branchwise-overlap-continuity}.
Then:
\begin{enumerate}[leftmargin=2em]
\item on every connected subsegment \(I'\subset I\) such that
  \[
  \nu_{\Lambda}>0
  \qquad
  \forall \Lambda\in I',
  \]
  the sign
  \begin{equation}
  \varepsilon_{\Lambda}
  :=
  \operatorname{sgn} m_I(\Lambda)
  =
  \operatorname{sgn}\!\Big(
  \overline{\rho}^{\mathrm{dom}}_{\Lambda}\big([\![\Xi^{(1)}_{\Lambda}]\!]\big)
  \Big)
  \in\{\pm1\}
  \label{eq:branchwise-epsilon}
  \end{equation}
  is constant on \(I'\);
\item if \(\nu_{\Lambda_0}>0\) at one point \(\Lambda_0\in I\), then there exists a neighborhood
  \(I_0\subset I\) of \(\Lambda_0\) on which \(\varepsilon_{\Lambda}\) is constant and equal to
  \(\varepsilon_{\Lambda_0}\);
\item under the present hypotheses, a change of sign can occur only at a cutoff where the
  descendant overlap closes
  \(
  \nu_{\Lambda}=0
  \).
\end{enumerate}
In particular, once dominant-line realization, quotient descent, and nonzero overlap are
established on an accepted branch segment and the projected descendant response varies continuously
there, the branchwise sign statement of \cref{thm:active-response-transport-minimal} follows.
\end{theorem}

\begin{proof}
By \cref{prop:branchwise-overlap-continuity}, \(m_I(\Lambda)\) is a continuous real-valued function
on \(I\).
On any connected subsegment \(I'\) where \(\nu_{\Lambda}=|m_I(\Lambda)|>0\), the function \(m_I\)
never vanishes.
A continuous nonvanishing real-valued function on a connected set has constant sign, proving item
(1).

Item (2) is the local version of the same argument: if \(m_I(\Lambda_0)\neq 0\), continuity gives a
neighborhood \(I_0\) on which \(m_I\) keeps the sign it has at \(\Lambda_0\).

For item (3), any sign change would force \(m_I\) to pass through zero by the intermediate value
theorem, hence \(\nu_{\Lambda}=|m_I(\Lambda)|=0\) at some intermediate cutoff.
\end{proof}

\begin{corollary}[Local orientation persistence from one-tick activity]
\label{cor:branchwise-orientation-local}
Let \(I\) be a connected accepted-branch segment and assume the hypotheses of
\cref{prop:branchwise-overlap-continuity}.
If the transported descendant class is dominant-active at one cutoff \(\Lambda_0\in I\), i.e.
\[
\nu_{\Lambda_0}>0,
\]
then there exists a neighborhood \(I_0\subset I\) of \(\Lambda_0\) such that
\[
\varepsilon_{\Lambda}=\varepsilon_{\Lambda_0}
\qquad
\forall \Lambda\in I_0.
\]
\end{corollary}

\begin{proof}
This is exactly item (2) of \cref{thm:branchwise-orientation-transport}.
\end{proof}

\subsection*{Witness-indexed seam realization and scalar bridge}
\label{subsec:witness-seam-bridge}

The bridge used below factors through a common LS witness datum rather than a direct identification
of a UOM seam source with a descendant representative.
On the UOM side one realizes a filler tag together with a detector witness by an admissible odd
welded source, measures the resulting dominant physical pairing, and then compares that scalar with
the odd representative attached to the same witness.

\begin{definition}[Accepted record payload of a welded odd source]
\label{def:accepted-record-payload}
Fix an accepted tick \(k\).
For an admissible welded odd trace-channel source \(\tau\), let
\[
\rho_{\Lambda_k}[\tau]
\]
denote the receiver state at cutoff \(\Lambda_k\) produced by that source in the welded UOM channel.
Define its accepted record payload by
\begin{equation}
\mathcal P^{\mathrm{acc}}_{\Lambda_k}[\tau]
:=
\Pi_{\Delta(\Lambda_k)}\!\big(\rho_{\Lambda_k}[\tau]\big).
\label{eq:accepted-record-payload}
\end{equation}
This is the diagonal/record part retained by the finite-cutoff acceptance map.
\end{definition}

\begin{definition}[Witness-indexed seam realization at an accepted tick]
\label{def:witness-seam-realization}
Fix an accepted tick \(k\).
Let \(y\) be a witness on the LS difference object, and write \(\mathrm{flip}(y)\) for the induced
LS involution.
A \emph{witness-indexed seam realization} at \(\Lambda_k\) is a choice of admissible welded odd
trace-channel sources
\[
\tau^{f,y}_{\Lambda_k},
\qquad
f\in\{\mathrm{diag},\mathrm{bridge}\},
\]
such that:
\begin{enumerate}[leftmargin=2em]
\item \textbf{filler toggle changes the odd seam sign:}
  \[
  \tau^{\mathrm{bridge},y}_{\Lambda_k}
  =
  -\,\tau^{\mathrm{diag},y}_{\Lambda_k};
  \]
\item \textbf{witness flip changes the odd seam sign:}
  \[
  \tau^{f,\mathrm{flip}(y)}_{\Lambda_k}
  =
  -\,\tau^{f,y}_{\Lambda_k}
  \qquad
  \forall f\in\{\mathrm{diag},\mathrm{bridge}\};
  \]
\item \textbf{accepted record payload is unchanged under filler toggle:}
  \[
  \mathcal P^{\mathrm{acc}}_{\Lambda_k}\!\big[\tau^{\mathrm{bridge},y}_{\Lambda_k}\big]
  =
  \mathcal P^{\mathrm{acc}}_{\Lambda_k}\!\big[\tau^{\mathrm{diag},y}_{\Lambda_k}\big].
  \]
\end{enumerate}
\end{definition}

\begin{remark}[Accepted-record invariance versus seam-scalar use]
\label{rem:accepted-record-vs-seam}
The accepted-record invariance clause in \cref{def:witness-seam-realization} encodes the physical
indistinguishability of filler toggle at the accepted-record layer.
The seam-scalar theorems below use only the induced odd sign relations in the source channel.
\end{remark}

\begin{definition}[Raw and normalized seam pairings]
\label{def:raw-seam-pairing}
Fix an accepted tick \(k\) and a witness-indexed seam realization as in
\cref{def:witness-seam-realization}.
For \(f\in\{\mathrm{diag},\mathrm{bridge}\}\) and witness \(y\), define the \emph{raw seam
pairing} by
\begin{equation}
\widetilde s_{\Lambda_k}(f,y)
:=
\big\langle\!\big\langle
\Pi_{\mathrm{phys},\Lambda_k}\mathsf T_{\Lambda_k}\!\big[\tau^{f,y}_{\Lambda_k}\big],\,
\delta C^{\mathrm{LoG}}_{\Lambda_k}
\big\rangle\!\big\rangle_{C_{\ast,\Lambda_k}}.
\label{eq:raw-seam-pairing}
\end{equation}
Whenever dominant-line realization holds so that
\(
\widehat{\delta C}^{\mathrm{LoG}}_{\Lambda_k}
\)
is defined, also set
\begin{equation}
s_{\Lambda_k}(f,y)
:=
\big\langle\!\big\langle
\Pi_{\mathrm{phys},\Lambda_k}\mathsf T_{\Lambda_k}\!\big[\tau^{f,y}_{\Lambda_k}\big],\,
\widehat{\delta C}^{\mathrm{LoG}}_{\Lambda_k}
\big\rangle\!\big\rangle_{C_{\ast,\Lambda_k}}.
\label{eq:normalized-seam-pairing}
\end{equation}
\end{definition}

\begin{proposition}[Filler and witness parity of the seam scalar]
\label{prop:seam-scalar-signs}
Fix an accepted tick \(k\) and assume a witness-indexed seam realization at \(\Lambda_k\).
Then the raw seam pairing satisfies
\begin{equation}
\widetilde s_{\Lambda_k}(\mathrm{bridge},y)
=
-\,\widetilde s_{\Lambda_k}(\mathrm{diag},y),
\qquad
\widetilde s_{\Lambda_k}(f,\mathrm{flip}(y))
=
-\,\widetilde s_{\Lambda_k}(f,y).
\label{eq:raw-seam-signs}
\end{equation}
If dominant-line realization also holds, then the same identities hold for the normalized seam
pairing:
\begin{equation}
s_{\Lambda_k}(\mathrm{bridge},y)
=
-\,s_{\Lambda_k}(\mathrm{diag},y),
\qquad
s_{\Lambda_k}(f,\mathrm{flip}(y))
=
-\,s_{\Lambda_k}(f,y).
\label{eq:normalized-seam-signs}
\end{equation}
\end{proposition}

\begin{proof}
The map
\(
\tau\mapsto \Pi_{\mathrm{phys},\Lambda_k}\mathsf T_{\Lambda_k}[\tau]
\)
is linear because \(\mathsf T_{\Lambda_k}\) and \(\Pi_{\mathrm{phys},\Lambda_k}\) are linear.
The Fisher pairing is bilinear.
Applying the sign relations of \cref{def:witness-seam-realization} therefore gives
\[
\Pi_{\mathrm{phys},\Lambda_k}\mathsf T_{\Lambda_k}\!\big[\tau^{\mathrm{bridge},y}_{\Lambda_k}\big]
=
-\,\Pi_{\mathrm{phys},\Lambda_k}\mathsf T_{\Lambda_k}\!\big[\tau^{\mathrm{diag},y}_{\Lambda_k}\big]
\]
and
\[
\Pi_{\mathrm{phys},\Lambda_k}\mathsf T_{\Lambda_k}\!\big[\tau^{f,\mathrm{flip}(y)}_{\Lambda_k}\big]
=
-\,\Pi_{\mathrm{phys},\Lambda_k}\mathsf T_{\Lambda_k}\!\big[\tau^{f,y}_{\Lambda_k}\big].
\]
Pairing with \(\delta C^{\mathrm{LoG}}_{\Lambda_k}\) proves \eqref{eq:raw-seam-signs}.
If dominant-line realization holds, then
\(
\widehat{\delta C}^{\mathrm{LoG}}_{\Lambda_k}
\)
is a fixed nonzero normalization of the same vector, so the same argument proves
\eqref{eq:normalized-seam-signs}.
\end{proof}

\begin{definition}[Detector-factorized seam coupling]
\label{def:detector-factorized-seam}
Fix an accepted tick \(k\).
A witness-indexed seam realization is called \emph{detector-factorized} if there exists a single
admissible welded odd trace-channel source
\[
\tau^{\mathrm{unit}}_{\Lambda_k}
\]
and the LS detector factor
\[
\overline{\mathrm{obs}}
:
\text{LS difference object}\to\mathbb R
\]
such that, for every filler tag \(f\) and witness \(y\),
\begin{equation}
\tau^{f,y}_{\Lambda_k}
=
\sigma_{\mathrm{fill}}(f)\,
\overline{\mathrm{obs}}(y)\,
\tau^{\mathrm{unit}}_{\Lambda_k},
\qquad
\sigma_{\mathrm{fill}}(\mathrm{diag})=+1,\quad
\sigma_{\mathrm{fill}}(\mathrm{bridge})=-1.
\label{eq:detector-factorized-seam}
\end{equation}
\end{definition}

\begin{theorem}[Seam faithfulness scalar]
\label{thm:seam-faithfulness-scalar}
Fix an accepted tick \(k\) and assume a detector-factorized witness-indexed seam realization at
\(\Lambda_k\).
Define the fixed-cutoff seam coefficient by
\begin{equation}
\widetilde c_{\Lambda_k}
:=
\big\langle\!\big\langle
\Pi_{\mathrm{phys},\Lambda_k}\mathsf T_{\Lambda_k}\!\big[\tau^{\mathrm{unit}}_{\Lambda_k}\big],\,
\delta C^{\mathrm{LoG}}_{\Lambda_k}
\big\rangle\!\big\rangle_{C_{\ast,\Lambda_k}}.
\label{eq:seam-faithfulness-coeff}
\end{equation}
Then, for every filler tag \(f\) and witness \(y\),
\begin{equation}
\widetilde s_{\Lambda_k}(f,y)
=
\sigma_{\mathrm{fill}}(f)\,
\widetilde c_{\Lambda_k}\,
\overline{\mathrm{obs}}(y).
\label{eq:raw-seam-factorization}
\end{equation}
In particular, if
\[
\widetilde c_{\Lambda_k}\neq 0
\qquad\text{and}\qquad
\overline{\mathrm{obs}}(y)\neq 0,
\]
then the raw seam pairing is nonzero:
\[
\widetilde s_{\Lambda_k}(f,y)\neq 0.
\]
Consequently the compressed welded LoG line is transversely physical at \(\Lambda_k\), hence
dominant-line realization holds at that cutoff by \cref{prop:dominant-line-equivalent}.
When dominant-line realization holds, the normalized seam coefficient
\begin{equation}
c_{\Lambda_k}
:=
\kappa_{\Lambda_k}^{-1/2}\,\widetilde c_{\Lambda_k}
\label{eq:normalized-seam-coeff}
\end{equation}
is defined and
\begin{equation}
s_{\Lambda_k}(f,y)
=
\sigma_{\mathrm{fill}}(f)\,
c_{\Lambda_k}\,
\overline{\mathrm{obs}}(y).
\label{eq:normalized-seam-factorization}
\end{equation}
\end{theorem}

\begin{proof}
Substitute the detector-factorized form \eqref{eq:detector-factorized-seam} into the definition
\eqref{eq:raw-seam-pairing}.
Linearity of \(\mathsf T_{\Lambda_k}\) and bilinearity of the Fisher pairing give
\[
\widetilde s_{\Lambda_k}(f,y)
=
\sigma_{\mathrm{fill}}(f)\,\overline{\mathrm{obs}}(y)\,
\big\langle\!\big\langle
\Pi_{\mathrm{phys},\Lambda_k}\mathsf T_{\Lambda_k}\!\big[\tau^{\mathrm{unit}}_{\Lambda_k}\big],\,
\delta C^{\mathrm{LoG}}_{\Lambda_k}
\big\rangle\!\big\rangle_{C_{\ast,\Lambda_k}},
\]
which is exactly \eqref{eq:raw-seam-factorization}.

If \(\widetilde c_{\Lambda_k}\neq 0\) and \(\overline{\mathrm{obs}}(y)\neq 0\), then
\(
\widetilde s_{\Lambda_k}(f,y)\neq 0
\).
Set
\[
B_{\Lambda_k}
:=
\Pi_{\mathrm{phys},\Lambda_k}\mathsf T_{\Lambda_k}\!\big[\tau^{f,y}_{\Lambda_k}\big].
\]
Then
\[
\big\langle\!\big\langle
\delta C^{\mathrm{LoG}}_{\Lambda_k},\,B_{\Lambda_k}
\big\rangle\!\big\rangle_{C_{\ast,\Lambda_k}}
=
\widetilde s_{\Lambda_k}(f,y)
\neq 0.
\]
By item (4) of \cref{prop:dominant-line-equivalent}, this proves transversality of the compressed
line and hence dominant-line realization at \(\Lambda_k\).

Once dominant-line realization holds, divide \eqref{eq:raw-seam-factorization} by
\(
\|\delta C^{\mathrm{LoG}}_{\Lambda_k}\|_{C_{\ast,\Lambda_k}}
=
\kappa_{\Lambda_k}^{1/2}
\)
to obtain \eqref{eq:normalized-seam-factorization}.
\end{proof}

\begin{remark}[Downstream role of the seam-coefficient hypothesis]
The fixed-cutoff nonvanishing hypothesis
\(
\widetilde c_{\Lambda_k}\neq 0
\)
is exactly the input that, once dominant-line realization defines
\(
c_{\Lambda_k}
\),
is imported downstream as the nondegeneracy clause
\(
c_{\Lambda_k}\neq 0
\)
inside the delayed selected-line package.
This note does not derive that clause from split-threshold analysis alone.
\end{remark}

\begin{definition}[Odd source representability on the characterization layer]
\label{def:odd-source-representability}
Fix an accepted tick \(k\).
Let \(\tau\) be an admissible welded odd trace-channel source.
We say that \(\tau\) is \emph{representable on the characterization layer} if there exists an odd
sufficiently local generator
\[
F_{\tau,\Lambda_k}\in E^-_{\Lambda_k}
\]
such that its intrinsic physical covariance response agrees with the UOM welded response:
\begin{equation}
\Pi_{\mathrm{phys},\Lambda_k}\mathcal R_{\Lambda_k}\!\big(F_{\tau,\Lambda_k}\big)
=
\Pi_{\mathrm{phys},\Lambda_k}\mathsf T_{\Lambda_k}[\tau].
\label{eq:odd-source-representability}
\end{equation}
This is the strong source-insertion realization statement.
\end{definition}

\begin{theorem}[UOM source-insertion representability theorem]
\label{thm:odd-source-representability}
Fix an accepted tick \(k\).
Let \(\tau\) be an admissible welded odd trace-channel source at cutoff \(\Lambda_k\), in the
finite-cutoff UOM sense: its receiver-head profile is band-limited at \(\Lambda_k\) and compactly
supported inside the accepted slab.
Then there exists a canonically associated odd sufficiently local generator
\[
F_{\tau,\Lambda_k}\in E^-_{\Lambda_k}
\]
such that
\begin{equation}
\Pi_{\mathrm{phys},\Lambda_k}\mathcal R_{\Lambda_k}\!\big(F_{\tau,\Lambda_k}\big)
=
\Pi_{\mathrm{phys},\Lambda_k}\mathsf T_{\Lambda_k}[\tau].
\label{eq:odd-source-representability-thm}
\end{equation}
In particular, every witness-indexed seam source used in the present note is representable on the
characterization layer.
\end{theorem}

\begin{proof}
By the finite-cutoff UOM Peierls package, admissible receiver-head sources at cutoff \(\Lambda_k\)
are band-limited boundary test data and the sufficiently local class
\(\mathcal F_{\rm s.loc}^{\Lambda_k}\) contains all compact smearings of local boundary densities.
Let
\[
F_{\tau,\Lambda_k}
\]
be the linear source functional obtained by pairing the odd trace-channel source profile \(\tau\)
with the corresponding odd trace-channel boundary field on the receiver head.
Because this is the linear case of the sufficiently local form and \(\tau\) is physically odd,
\[
F_{\tau,\Lambda_k}\in E^-_{\Lambda_k}.
\]

Now apply the UOM source insertion map from the finite-cutoff retarded-kernel package.
The source \(\tau\) is sent to the bulk forcing term in the linearized equation; the unique
retarded solution is obtained by the retarded Green operator; and boundary restriction gives the
retarded covariance variation. By the UOM identification of the induced retarded kernel with the SK
mixed second derivative, this covariance variation is exactly the welded source-to-covariance
response
\[
\mathsf T_{\Lambda_k}[\tau].
\]

On the other hand, by the definition of the intrinsic covariance-response map, the matrix elements
of
\(
\mathcal R_{\Lambda_k}(F_{\tau,\Lambda_k})
\)
are computed by the Peierls action of \(F_{\tau,\Lambda_k}\) on the sufficiently local quadratic
probes representing boundary covariance entries.
That Peierls variation is precisely the same retarded second-moment variation generated by the
source insertion \(\tau\).
Therefore
\[
\mathcal R_{\Lambda_k}(F_{\tau,\Lambda_k})
=
\mathsf T_{\Lambda_k}[\tau]
\]
as covariance tangents.
Applying the physical projector \(\Pi_{\mathrm{phys},\Lambda_k}\) gives
\eqref{eq:odd-source-representability-thm}.

The witness-indexed seam sources of \cref{def:witness-seam-realization} are, by construction, live
admissible welded odd trace-channel sources at the accepted tick, so the same argument applies to
each of them.
\end{proof}

\begin{proposition}[Dominant pairing for a representable odd source]
\label{prop:odd-source-pairing}
Fix an accepted tick \(k\) and let \(\tau\) be an admissible welded odd trace-channel source that is
representable in the sense of \cref{thm:odd-source-representability}.
Then
\begin{equation}
\big\langle\!\big\langle
\Pi_{\mathrm{phys},\Lambda_k}\mathcal R_{\Lambda_k}\!\big(F_{\tau,\Lambda_k}\big),\,
\delta C^{\mathrm{LoG}}_{\Lambda_k}
\big\rangle\!\big\rangle_{C_{\ast,\Lambda_k}}
=
\big\langle\!\big\langle
\Pi_{\mathrm{phys},\Lambda_k}\mathsf T_{\Lambda_k}[\tau],\,
\delta C^{\mathrm{LoG}}_{\Lambda_k}
\big\rangle\!\big\rangle_{C_{\ast,\Lambda_k}}.
\label{eq:representable-raw-pairing}
\end{equation}
If dominant-line realization also holds, then
\begin{equation}
\rho^{\mathrm{dom}}_{\Lambda_k}\!\big(F_{\tau,\Lambda_k}\big)
=
\big\langle\!\big\langle
\Pi_{\mathrm{phys},\Lambda_k}\mathsf T_{\Lambda_k}[\tau],\,
\widehat{\delta C}^{\mathrm{LoG}}_{\Lambda_k}
\big\rangle\!\big\rangle_{C_{\ast,\Lambda_k}}.
\label{eq:representable-normalized-pairing}
\end{equation}
\end{proposition}

\begin{proof}
Equation \eqref{eq:odd-source-representability} identifies the two physical covariance tangents.
Pairing both sides with \(\delta C^{\mathrm{LoG}}_{\Lambda_k}\) gives
\eqref{eq:representable-raw-pairing}.
If dominant-line realization holds, pair with the normalized vector
\(
\widehat{\delta C}^{\mathrm{LoG}}_{\Lambda_k}
\)
instead and use the definition \eqref{eq:rho-dom-def} of
\(
\rho^{\mathrm{dom}}_{\Lambda_k}
\)
to obtain \eqref{eq:representable-normalized-pairing}.
\end{proof}

\begin{corollary}[Unit welded odd source representability and scalar identification]
\label{cor:unit-source-representability}
Fix an accepted tick \(k\).
Assume a detector-factorized witness-indexed seam realization at \(\Lambda_k\) in the sense of
\cref{def:detector-factorized-seam}, and assume dominant-line realization at \(\Lambda_k\).
Then the unit welded odd source
\(
\tau^{\mathrm{unit}}_{\Lambda_k}
\)
is representable on the characterization layer, with canonical representable generator
\[
F_{\tau^{\mathrm{unit}}_{\Lambda_k},\Lambda_k}\in E^-_{\Lambda_k},
\]
and
\begin{equation}
\rho^{\mathrm{dom}}_{\Lambda_k}
\!\bigl(
F_{\tau^{\mathrm{unit}}_{\Lambda_k},\Lambda_k}
\bigr)
=
\big\langle\!\big\langle
\Pi_{\mathrm{phys},\Lambda_k}\mathsf T_{\Lambda_k}\!\big[\tau^{\mathrm{unit}}_{\Lambda_k}\big],\,
\widehat{\delta C}^{\mathrm{LoG}}_{\Lambda_k}
\big\rangle\!\big\rangle_{C_{\ast,\Lambda_k}}
=
c_{\Lambda_k}.
\label{eq:unit-source-representability-scalar}
\end{equation}
\end{corollary}

\begin{proof}
Representability of
\(
\tau^{\mathrm{unit}}_{\Lambda_k}
\)
is the specialization of \cref{thm:odd-source-representability} to the admissible welded odd trace
source singled out by \cref{def:detector-factorized-seam}.
Because dominant-line realization holds, \cref{prop:odd-source-pairing} applies to that
representable source and gives the first equality in \eqref{eq:unit-source-representability-scalar}.
The second equality is exactly the definition of the normalized seam coefficient
\(
c_{\Lambda_k}
\)
from \cref{thm:seam-faithfulness-scalar}.
\end{proof}

\begin{definition}[Scalar comparison shadow of the witness-indexed bridge]
\label{def:witness-descendant-comparison}
Fix an accepted tick \(k\).
Let \(y\) be an LS witness and let
\(
\tau^{f,y}_{\Lambda_k}
\)
be a witness-indexed seam realization source for a filler tag \(f\).
A \emph{scalar comparison shadow} for that data is a choice of representative
\[
J^{f,y}_{\Lambda_k}\in[\Xi^{(1)}_{\Lambda_k}]
\]
such that the dominant raw pairing of the 2T-side representative agrees with the dominant raw
pairing of the corresponding representable UOM source:
\begin{equation}
\big\langle\!\big\langle
\Pi_{\mathrm{phys},\Lambda_k}\mathcal R_{\Lambda_k}\!\big(J^{f,y}_{\Lambda_k}\big),\,
\delta C^{\mathrm{LoG}}_{\Lambda_k}
\big\rangle\!\big\rangle_{C_{\ast,\Lambda_k}}
=
\big\langle\!\big\langle
\Pi_{\mathrm{phys},\Lambda_k}\mathcal R_{\Lambda_k}\!\big(F_{\tau^{f,y}_{\Lambda_k},\Lambda_k}\big),\,
\delta C^{\mathrm{LoG}}_{\Lambda_k}
\big\rangle\!\big\rangle_{C_{\ast,\Lambda_k}}.
\label{eq:witness-descendant-comparison}
\end{equation}
When dominant-line realization holds, this is equivalently the equality of normalized dominant
pairings
\begin{equation}
\rho^{\mathrm{dom}}_{\Lambda_k}\!\big(J^{f,y}_{\Lambda_k}\big)
=
\rho^{\mathrm{dom}}_{\Lambda_k}\!\big(F_{\tau^{f,y}_{\Lambda_k},\Lambda_k}\big).
\label{eq:witness-descendant-comparison-normalized}
\end{equation}
\end{definition}

\begin{definition}[Witness-indexed source-compatible descendant realization]
\label{def:source-compatible-witness-realization}
Fix an accepted tick \(k\), a filler tag \(f\), a witness \(y\), and a witness-indexed seam source
\(
\tau^{f,y}_{\Lambda_k}
\).
A representative
\[
J^{f,y}_{\Lambda_k}\in[\Xi^{(1)}_{\Lambda_k}]
\]
is called a \emph{witness-indexed source-compatible descendant realization} if it realizes, on the
same LS witness datum, the Limited Scope output and odd Hamiltonization clauses from item
\textbf{(M1)} and item \textbf{(M4)} of Definition~\ref{main-def:uom-main-theorem-assumptions} and differs
from the source-insertion representative only by an ambiguity direction:
\begin{equation}
J^{f,y}_{\Lambda_k}
-
F_{\tau^{f,y}_{\Lambda_k},\Lambda_k}
\in
\mathcal A^-_{\Lambda_k},
\label{eq:source-compatible-witness-realization}
\end{equation}
where
\(
F_{\tau^{f,y}_{\Lambda_k},\Lambda_k}
\)
is the representable source generator furnished by
\cref{thm:odd-source-representability}.
Equivalently, the witness-indexed transported descendant class and the witness-indexed UOM
source-insertion realization define the same point of the ambiguity quotient
\(
E^-_{\Lambda_k}/\mathcal A^-_{\Lambda_k}
\).
\end{definition}
\begin{theorem}[Same-witness descended-scalar comparison from a source-compatible descendant realization]
\label{thm:witness-descended-scalar-comparison}
Fix an accepted tick \(k\), assume dominant-line realization at \(\Lambda_k\), and let
\(
\tau^{f,y}_{\Lambda_k}
\)
be a witness-indexed seam realization source.
Assume:
\begin{enumerate}[leftmargin=2em]
\item the source \(\tau^{f,y}_{\Lambda_k}\) is represented by the canonical generator
  \(
  F_{\tau^{f,y}_{\Lambda_k},\Lambda_k}
  \)
  from \cref{thm:odd-source-representability};
\item there exists a witness-indexed source-compatible descendant realization
  \(
  J^{f,y}_{\Lambda_k}\in[\Xi^{(1)}_{\Lambda_k}]
  \)
  in the sense of \cref{def:source-compatible-witness-realization};
\item the three vanishing statements of \cref{prop:dominant-quotient-generators} hold at
  \(\Lambda_k\).
\end{enumerate}
Then the raw and normalized comparison identities of
\cref{def:witness-descendant-comparison} hold:
\begin{align}
\big\langle\!\big\langle
\Pi_{\mathrm{phys},\Lambda_k}\mathcal R_{\Lambda_k}\!\big(J^{f,y}_{\Lambda_k}\big),\,
\delta C^{\mathrm{LoG}}_{\Lambda_k}
\big\rangle\!\big\rangle_{C_{\ast,\Lambda_k}}
&=
\big\langle\!\big\langle
\Pi_{\mathrm{phys},\Lambda_k}\mathcal R_{\Lambda_k}\!\big(F_{\tau^{f,y}_{\Lambda_k},\Lambda_k}\big),\,
\delta C^{\mathrm{LoG}}_{\Lambda_k}
\big\rangle\!\big\rangle_{C_{\ast,\Lambda_k}},
\label{eq:witness-descendant-comparison-thm-raw}
\\
\rho^{\mathrm{dom}}_{\Lambda_k}\!\big(J^{f,y}_{\Lambda_k}\big)
&=
\rho^{\mathrm{dom}}_{\Lambda_k}\!\big(F_{\tau^{f,y}_{\Lambda_k},\Lambda_k}\big).
\label{eq:witness-descendant-comparison-thm-normalized}
\end{align}
Moreover,
\begin{equation}
\rho^{\mathrm{dom}}_{\Lambda_k}\!\big(F_{\tau^{f,y}_{\Lambda_k},\Lambda_k}\big)
=
\big\langle\!\big\langle
\Pi_{\mathrm{phys},\Lambda_k}\mathsf T_{\Lambda_k}\!\big[\tau^{f,y}_{\Lambda_k}\big],\,
\widehat{\delta C}^{\mathrm{LoG}}_{\Lambda_k}
\big\rangle\!\big\rangle_{C_{\ast,\Lambda_k}}.
\label{eq:witness-descendant-comparison-thm-seam}
\end{equation}
If quotient descent also holds at \(\Lambda_k\), then the common value equals the descended
dominant scalar on the transported descendant affine class:
\begin{equation}
\overline{\rho}^{\mathrm{dom}}_{\Lambda_k}\big([\![\Xi^{(1)}_{\Lambda_k}]\!]\big)
=
\rho^{\mathrm{dom}}_{\Lambda_k}\!\big(J^{f,y}_{\Lambda_k}\big)
=
\rho^{\mathrm{dom}}_{\Lambda_k}\!\big(F_{\tau^{f,y}_{\Lambda_k},\Lambda_k}\big).
\label{eq:witness-descendant-comparison-thm-descended}
\end{equation}
\end{theorem}

\begin{proof}
Set
\[
A^{f,y}_{\Lambda_k}
:=
J^{f,y}_{\Lambda_k}
-
F_{\tau^{f,y}_{\Lambda_k},\Lambda_k}.
\]
By source compatibility,
\(
A^{f,y}_{\Lambda_k}\in\mathcal A^-_{\Lambda_k}
\).

Because dominant-line realization holds and the three vanishing statements of
\cref{prop:dominant-quotient-generators} are assumed, that proposition gives
\[
\mathcal A^-_{\Lambda_k}\subset\ker\rho^{\mathrm{dom}}_{\Lambda_k}.
\]
Hence
\[
\rho^{\mathrm{dom}}_{\Lambda_k}\!\big(J^{f,y}_{\Lambda_k}\big)
=
\rho^{\mathrm{dom}}_{\Lambda_k}\!\big(
F_{\tau^{f,y}_{\Lambda_k},\Lambda_k}
+
A^{f,y}_{\Lambda_k}
\big)
=
\rho^{\mathrm{dom}}_{\Lambda_k}\!\big(F_{\tau^{f,y}_{\Lambda_k},\Lambda_k}\big),
\]
proving \eqref{eq:witness-descendant-comparison-thm-normalized}.

Since dominant-line realization implies
\(
\delta C^{\mathrm{LoG}}_{\Lambda_k}\neq 0
\)
and
\(
\rho^{\mathrm{dom}}_{\Lambda_k}
=
\kappa_{\Lambda_k}^{-1/2}\widetilde{\rho}^{\mathrm{dom}}_{\Lambda_k},
\)
the raw pairing equality \eqref{eq:witness-descendant-comparison-thm-raw} is just the same identity
multiplied by the positive factor \(\kappa_{\Lambda_k}^{1/2}\).

Equation \eqref{eq:witness-descendant-comparison-thm-seam} is exactly the normalized pairing formula
of \cref{prop:odd-source-pairing} applied to the representable source
\(
\tau^{f,y}_{\Lambda_k}
\).

If quotient descent also holds, then \cref{cor:dominant-quotient-constancy} gives
\[
\rho^{\mathrm{dom}}_{\Lambda_k}\!\big(J^{f,y}_{\Lambda_k}\big)
=
\overline{\rho}^{\mathrm{dom}}_{\Lambda_k}\big([\![\Xi^{(1)}_{\Lambda_k}]\!]\big)
\]
for every representative
\(
J^{f,y}_{\Lambda_k}\in[\Xi^{(1)}_{\Lambda_k}],
\)
which yields \eqref{eq:witness-descendant-comparison-thm-descended} after combining with
\eqref{eq:witness-descendant-comparison-thm-normalized}.
\end{proof}

\begin{theorem}[Fixed-cutoff seam witness criterion for descendant overlap]
\label{thm:fixed-cutoff-seam-witness}
Fix an accepted tick \(k\).
Assume:
\begin{enumerate}[leftmargin=2em]
\item a detector-factorized witness-indexed seam realization at \(\Lambda_k\);
\item a witness \(y\) with
  \[
  \overline{\mathrm{obs}}(y)\neq 0;
  \]
\item the seam coefficient is nonzero,
  \[
  \widetilde c_{\Lambda_k}\neq 0;
  \]
\item for the chosen filler tag \(f\), the source \(\tau^{f,y}_{\Lambda_k}\) is representable on the
  characterization layer in the sense of \cref{def:odd-source-representability};
\item there exists a witness-indexed source-compatible descendant realization in the sense of
  \cref{def:source-compatible-witness-realization}.
\end{enumerate}
Then:
\begin{enumerate}[leftmargin=2em]
\item the compressed welded LoG line is transversely physical at \(\Lambda_k\), hence dominant-line
  realization holds at that cutoff;
\item the chosen 2T representative satisfies
  \[
  \rho^{\mathrm{dom}}_{\Lambda_k}\!\big(J^{f,y}_{\Lambda_k}\big)\neq 0;
  \]
\item therefore the transported descendant class has nonzero overlap with the canonically selected
  compressed dominant odd line at \(\Lambda_k\);
\item if quotient descent also holds at \(\Lambda_k\), then
  \cref{thm:descendant-overlap} applies and the transported descendant class is dominant-active with
  well-defined sign \(\varepsilon_{\Lambda_k}\).
\end{enumerate}
\end{theorem}

\begin{proof}
By \cref{thm:seam-faithfulness-scalar},
\[
\widetilde s_{\Lambda_k}(f,y)
=
\sigma_{\mathrm{fill}}(f)\,
\widetilde c_{\Lambda_k}\,
\overline{\mathrm{obs}}(y)
\neq 0.
\]
The same theorem therefore gives item (1): the compressed line is transversely physical, hence
dominant-line realization holds at \(\Lambda_k\).

By \cref{prop:odd-source-pairing},
\[
\big\langle\!\big\langle
\Pi_{\mathrm{phys},\Lambda_k}\mathcal R_{\Lambda_k}\!\big(F_{\tau^{f,y}_{\Lambda_k},\Lambda_k}\big),\,
\delta C^{\mathrm{LoG}}_{\Lambda_k}
\big\rangle\!\big\rangle_{C_{\ast,\Lambda_k}}
=
\widetilde s_{\Lambda_k}(f,y)
\neq 0.
\]
Applying the comparison identity \eqref{eq:witness-descendant-comparison} now gives
\[
\big\langle\!\big\langle
\Pi_{\mathrm{phys},\Lambda_k}\mathcal R_{\Lambda_k}\!\big(J^{f,y}_{\Lambda_k}\big),\,
\delta C^{\mathrm{LoG}}_{\Lambda_k}
\big\rangle\!\big\rangle_{C_{\ast,\Lambda_k}}
\neq 0.
\]
Because dominant-line realization has already been established, \(\delta C^{\mathrm{LoG}}_{\Lambda_k}\)
is nonzero and \(\widehat{\delta C}^{\mathrm{LoG}}_{\Lambda_k}\) is defined.
By \cref{prop:descendant-overlap-witness}, this is equivalent to
\(
\rho^{\mathrm{dom}}_{\Lambda_k}(J^{f,y}_{\Lambda_k})\neq 0
\),
proving item (2).

Item (3) is exactly the nonzero-overlap condition from
\cref{prop:descendant-overlap-equivalent}.
If quotient descent is also available, apply \cref{thm:descendant-overlap} to obtain item (4).
\end{proof}

\subsection*{Smooth-taper branch continuity and witness transport}
\label{subsec:smooth-taper-branch-continuity}

The previous fixed-cutoff theorem isolates the nonzero witness needed for activity.
To transport that witness along an accepted branch one needs a continuity statement for the welded
response.
At theorem level, the clean analytic choice is the smooth receiver taper.
The hard projector may still be used, but only on branch segments where its rank is constant.

\begin{lemma}[Smooth-taper continuity of the receiver band projector]
\label{lem:smooth-taper-projector}
Let \(I\) be a connected accepted-branch segment.
Suppose the receiver band is implemented by a smooth taper
\[
P^{\mathrm{(rec)}}_{\mathrm{band}}(\Lambda)
=
\sum_{\ell,m} w_\ell(\Lambda)\,|Y_{\ell m}\rangle\langle Y_{\ell m}|,
\qquad
0\le w_\ell(\Lambda)\le 1,
\]
on the receiver one-particle space, where
\[
\sup_{\ell\ge 0}|w_\ell(\Lambda')-w_\ell(\Lambda)|
\longrightarrow 0
\qquad
\text{as }\Lambda'\to\Lambda
\]
for every \(\Lambda\in I\).
Then
\[
\Lambda\longmapsto P^{\mathrm{(rec)}}_{\mathrm{band}}(\Lambda)
\]
is continuous on \(I\) in operator norm.
If instead the hard projector
\(
\sum_{\ell\le \ell_\star(\Lambda)}|Y_{\ell m}\rangle\langle Y_{\ell m}|
\)
is used, the same conclusion holds on every connected subsegment of \(I\) on which
\(\ell_\star(\Lambda)\) is constant.
\end{lemma}

\begin{proof}
In the spherical-harmonic basis \(\{Y_{\ell m}\}\), the difference
\[
P^{\mathrm{(rec)}}_{\mathrm{band}}(\Lambda')
-
P^{\mathrm{(rec)}}_{\mathrm{band}}(\Lambda)
\]
is diagonal with eigenvalues \(w_\ell(\Lambda')-w_\ell(\Lambda)\).
Therefore
\[
\big\|
P^{\mathrm{(rec)}}_{\mathrm{band}}(\Lambda')
-
P^{\mathrm{(rec)}}_{\mathrm{band}}(\Lambda)
\big\|_{\mathrm{op}}
=
\sup_{\ell\ge 0}|w_\ell(\Lambda')-w_\ell(\Lambda)|.
\]
The assumed uniform continuity of the taper coefficients is exactly the statement that this operator
norm tends to zero as \(\Lambda'\to\Lambda\).

For the hard projector, the coefficients take only the values \(0\) and \(1\), so discontinuities can
occur only when \(\ell_\star(\Lambda)\) jumps.
On a subsegment where \(\ell_\star(\Lambda)\) is constant, the projector is constant and therefore
continuous.
\end{proof}

\begin{theorem}[Smooth-taper branch continuity of the welded physical response]
\label{thm:smooth-taper-welded-response}
Let \(I\) be a connected accepted-branch segment.
Write the welded source-to-covariance map in the UOM factorized form
\begin{equation}
\mathsf T_\Lambda
=
\Pi_{L(\Lambda)}M_{\varepsilon,\Lambda}
\circ
\mathsf K_{{\rm comp},\Lambda}
\circ
\mathsf G_{{\rm ret},\Lambda}
\circ
\mathsf W_{{\rm src},\Lambda}.
\label{eq:factorized-welded-map}
\end{equation}
Assume:
\begin{enumerate}[leftmargin=2em]
\item the receiver band is implemented by a smooth taper satisfying
  \cref{lem:smooth-taper-projector};
\item the operator families
  \[
  \Lambda\longmapsto
  \Pi_{L(\Lambda)}M_{\varepsilon,\Lambda},\qquad
  \Lambda\longmapsto
  \mathsf K_{{\rm comp},\Lambda},\qquad
  \Lambda\longmapsto
  \mathsf G_{{\rm ret},\Lambda},\qquad
  \Lambda\longmapsto
  \mathsf W_{{\rm src},\Lambda},\qquad
  \Lambda\longmapsto
  \Pi_{\mathrm{phys},\Lambda}
  \]
  are continuous on \(I\) in operator norm on the corresponding band-limited source and covariance
  tangent spaces;
\item the source family
  \[
  \Lambda\longmapsto \tau_\Lambda
  \]
  is continuous on \(I\) in the welded source norm.
\end{enumerate}
Then the physical covariance response
\begin{equation}
X_\Lambda[\tau_\Lambda]
:=
\Pi_{\mathrm{phys},\Lambda}\mathsf T_\Lambda[\tau_\Lambda]
\label{eq:XLambda-tau-def}
\end{equation}
depends continuously on \(\Lambda\in I\) in the Fisher norm.
If the hard receiver projector is used instead, the same conclusion holds on each connected
subsegment where \(\ell_\star(\Lambda)\) is constant.
\end{theorem}

\begin{proof}
By \cref{lem:smooth-taper-projector}, the first factor
\(
\Pi_{L(\Lambda)}M_{\varepsilon,\Lambda}
\)
is operator-norm continuous along the branch.
By assumption, the remaining three UOM factors in \eqref{eq:factorized-welded-map} and the physical
projector \(\Pi_{\mathrm{phys},\Lambda}\) are also operator-norm continuous.
Hence the composed operator family
\[
\Lambda\longmapsto
\Pi_{\mathrm{phys},\Lambda}\mathsf T_\Lambda
\]
is continuous in operator norm on \(I\).
Applying this continuous family to the continuous source vector \(\tau_\Lambda\) gives continuity of
\(
X_\Lambda[\tau_\Lambda]
\)
in the Fisher norm.

For the hard projector, the same argument applies on every connected subsegment where
\(\ell_\star(\Lambda)\) is constant, because there the first factor is constant.
\end{proof}

\begin{corollary}[Branchwise continuity of the seam coefficient]
\label{cor:branchwise-seam-coeff}
Let \(I\) be a connected accepted-branch segment and assume:
\begin{enumerate}[leftmargin=2em]
\item the hypotheses of \cref{thm:dominant-line-persistence} on \(I\);
\item a detector-factorized branchwise seam realization on \(I\) with continuous unit source
  \[
  \Lambda\longmapsto \tau^{\mathrm{unit}}_{\Lambda};
  \]
\item the hypotheses of \cref{thm:smooth-taper-welded-response} for that unit source.
\end{enumerate}
Define the branchwise normalized seam coefficient by
\begin{equation}
c_\Lambda
:=
\big\langle\!\big\langle
\Pi_{\mathrm{phys},\Lambda}\mathsf T_\Lambda\!\big[\tau^{\mathrm{unit}}_{\Lambda}\big],\,
\widehat{\delta C}^{\mathrm{LoG}}_{\Lambda}
\big\rangle\!\big\rangle_{C_{\ast,\Lambda}}.
\label{eq:branchwise-seam-coeff}
\end{equation}
Then \(c_\Lambda\) is continuous on \(I\).
Consequently, for every fixed filler tag \(f\) and witness \(y\), the normalized seam scalar
\[
s_\Lambda(f,y)
=
\sigma_{\mathrm{fill}}(f)\,
c_\Lambda\,
\overline{\mathrm{obs}}(y)
\]
is continuous on \(I\).
\end{corollary}

\begin{proof}
By \cref{thm:smooth-taper-welded-response}, the response vector
\[
\Lambda\longmapsto
\Pi_{\mathrm{phys},\Lambda}\mathsf T_\Lambda\!\big[\tau^{\mathrm{unit}}_{\Lambda}\big]
\]
is continuous in the Fisher norm.
By \cref{thm:dominant-line-persistence}, the normalized dominant vector
\(
\widehat{\delta C}^{\mathrm{LoG}}_{\Lambda}
\)
is also continuous on \(I\).
Pairing these two continuous vectors in the Fisher metric gives continuity of \(c_\Lambda\).
The final statement is then immediate from
\eqref{eq:normalized-seam-factorization}.
\end{proof}

\begin{theorem}[Branchwise seam witness transport]
\label{thm:branchwise-seam-witness}
Let \(I\) be a connected accepted-branch segment.
Assume:
\begin{enumerate}[leftmargin=2em]
\item the hypotheses of \cref{cor:branchwise-seam-coeff};
\item quotient descent holds for every \(\Lambda\in I\);
\item there exists a fixed LS witness \(y\) with
  \[
  \overline{\mathrm{obs}}(y)\neq 0;
  \]
\item for some fixed filler tag \(f\), there is a branchwise family of representatives
  \[
  J^{f,y}_{\Lambda}\in[\Xi^{(1)}_{\Lambda}]
  \qquad
  (\Lambda\in I)
  \]
  such that, for every \(\Lambda\in I\), the corresponding seam source
  \(\tau^{f,y}_{\Lambda}\) is representable on the characterization layer and the 2T-side
  witness-to-descendant comparison identity
  \eqref{eq:witness-descendant-comparison-normalized} holds.
\end{enumerate}
Then:
\begin{enumerate}[leftmargin=2em]
\item the branchwise dominant-response scalar of the chosen descendant representatives is
  \begin{equation}
  \rho^{\mathrm{dom}}_{\Lambda}\!\big(J^{f,y}_{\Lambda}\big)
  =
  \sigma_{\mathrm{fill}}(f)\,
  c_\Lambda\,
  \overline{\mathrm{obs}}(y);
  \label{eq:branchwise-seam-descendant-scalar}
  \end{equation}
\item the scalar function
  \(
  \Lambda\mapsto \rho^{\mathrm{dom}}_{\Lambda}(J^{f,y}_{\Lambda})
  \)
  is continuous on \(I\);
\item on every connected subsegment \(I'\subset I\) on which \(c_\Lambda\neq 0\), the sign
  \[
  \operatorname{sgn}\!\big(
  \rho^{\mathrm{dom}}_{\Lambda}(J^{f,y}_{\Lambda})
  \big)
  \]
  is constant;
\item because quotient descent holds, the same constant sign is the transported orientation sign of
  the descendant affine class on \(I'\).
\end{enumerate}
\end{theorem}

\begin{proof}
For each \(\Lambda\in I\), representability of \(\tau^{f,y}_{\Lambda}\) and
\cref{prop:odd-source-pairing} give
\[
\rho^{\mathrm{dom}}_{\Lambda}\!\big(F_{\tau^{f,y}_{\Lambda},\Lambda}\big)
=
\big\langle\!\big\langle
\Pi_{\mathrm{phys},\Lambda}\mathsf T_\Lambda\!\big[\tau^{f,y}_{\Lambda}\big],\,
\widehat{\delta C}^{\mathrm{LoG}}_{\Lambda}
\big\rangle\!\big\rangle_{C_{\ast,\Lambda}}
=
s_\Lambda(f,y).
\]
Applying the branchwise comparison identity
\eqref{eq:witness-descendant-comparison-normalized} then yields
\[
\rho^{\mathrm{dom}}_{\Lambda}\!\big(J^{f,y}_{\Lambda}\big)
=
s_\Lambda(f,y).
\]
Now use \cref{cor:branchwise-seam-coeff} to substitute
\(
s_\Lambda(f,y)
=
\sigma_{\mathrm{fill}}(f)c_\Lambda\overline{\mathrm{obs}}(y)
\),
which proves item (1).
Item (2) follows because \(c_\Lambda\) is continuous and
\(\overline{\mathrm{obs}}(y)\) is a fixed nonzero scalar.
On any connected subsegment where \(c_\Lambda\neq 0\), the right-hand side of
\eqref{eq:branchwise-seam-descendant-scalar} is a continuous nonvanishing real-valued function, so
its sign is constant.
This proves item (3).

Finally, quotient descent implies that \(\rho^{\mathrm{dom}}_{\Lambda}\) is constant on
\([\Xi^{(1)}_{\Lambda}]\), so the sign of one representative is the sign of the whole descendant
class.
This gives item (4).
\end{proof}

\begin{corollary}[Seam package discharge of the minimal scalar transport input]
\label{cor:branchwise-seam-minimal-input}
Let \(I\) be a connected accepted-branch segment and assume the hypotheses of
\cref{thm:branchwise-seam-witness}.
Then the branchwise family
\[
J^{f,y}_{\Lambda}\in[\Xi^{(1)}_{\Lambda}]
\qquad
(\Lambda\in I)
\]
satisfies
\[
\Lambda\longmapsto
\rho^{\mathrm{dom}}_{\Lambda}\!\big(J^{f,y}_{\Lambda}\big)
=
\sigma_{\mathrm{fill}}(f)\,
c_\Lambda\,
\overline{\mathrm{obs}}(y)
\]
continuously on \(I\).
On every connected subsegment \(I'\subset I\) such that
\[
c_\Lambda\neq 0
\qquad
\forall \Lambda\in I',
\]
assumption \textup{(3)} of \cref{thm:active-response-transport-minimal} holds in its minimal
scalar form with
\[
J^{\mathrm{desc}}_{\Lambda}:=J^{f,y}_{\Lambda}
\qquad
(\Lambda\in I').
\]
If, in addition, the three generator-wise vanishing statements of
\cref{prop:dominant-quotient-generators} hold for every \(\Lambda\in I'\), then
\cref{thm:active-response-transport-minimal} applies on \(I'\).
\end{corollary}

\begin{proof}
Continuity of the displayed scalar function is exactly item (2) of
\cref{thm:branchwise-seam-witness}, together with the explicit formula of item (1).
On \(I'\), the hypotheses
\(
c_\Lambda\neq 0
\)
and
\(
\overline{\mathrm{obs}}(y)\neq 0
\)
imply
\[
\rho^{\mathrm{dom}}_{\Lambda}\!\big(J^{f,y}_{\Lambda}\big)\neq 0
\qquad
\forall \Lambda\in I'.
\]
Together with the continuity statement already proved above, this is exactly assumption
\textup{(3)} of \cref{thm:active-response-transport-minimal} in its minimal scalar form.
The hypotheses of \cref{thm:branchwise-seam-witness} already include those of
\cref{cor:branchwise-seam-coeff}, and hence the dominant-line persistence input of
\cref{thm:active-response-transport-minimal}.
If the three generator-wise vanishing statements also hold on \(I'\), then every hypothesis of
\cref{thm:active-response-transport-minimal} is in force there.
\end{proof}

\begin{theorem}[Finite-cutoff Active Response Transport Theorem]
\label{thm:active-response-transport-minimal}
Let \(I\) be a connected accepted-branch segment.
For each \(\Lambda\in I\), let
\[
[\Xi^{(1)}_{\Lambda}]
:=
J_{0,\Lambda}+\mathcal A^-_{\Lambda}
\subset E^-_{\Lambda}
\]
be the parity-transported odd descendant affine class modulo the ambiguity space of
\cref{def:transport-ambiguity-subspace}.
Assume:
\begin{enumerate}[leftmargin=2em]
\item the hypotheses of \cref{thm:dominant-line-persistence} hold on \(I\);
\item for every \(\Lambda\in I\), the three vanishing statements of
  \cref{prop:dominant-quotient-generators} hold;
\item there exists a branchwise family of representatives
  \[
  J^{\mathrm{desc}}_{\Lambda}\in[\Xi^{(1)}_{\Lambda}]
  \qquad
  (\Lambda\in I)
  \]
  such that
  \[
  \rho^{\mathrm{dom}}_{\Lambda}(J^{\mathrm{desc}}_{\Lambda})
  \neq 0
  \qquad
  \forall \Lambda\in I,
  \]
  and such that the dominant-response scalar
  \[
  \Lambda\longmapsto
  \rho^{\mathrm{dom}}_{\Lambda}(J^{\mathrm{desc}}_{\Lambda})
  \]
  is continuous on \(I\).
\end{enumerate}
Then, for every \(\Lambda\in I\):
\begin{enumerate}[leftmargin=2em]
\item the compressed dominant covariance tangent is nonzero,
  \[
  \delta C^{\mathrm{LoG}}_{\Lambda}\neq 0,
  \]
  so the dominant line \(L^{\mathrm{dom}}_{\Lambda}\) of
  \cref{def:dominant-response-functional} is canonically defined;
\item the dominant-response coefficient \(\rho^{\mathrm{dom}}_{\Lambda}\) annihilates the odd
  ambiguity space,
  \[
  \rho^{\mathrm{dom}}_{\Lambda}(A)=0
  \qquad
  \forall A\in\mathcal A^-_{\Lambda};
  \]
\item the transported descendant affine class carries a unique nonzero dominant-response value
  \[
  \rho^{\mathrm{dom}}_{\Lambda}(J)
  =
  \overline{\rho}^{\mathrm{dom}}_{\Lambda}\big([\![\Xi^{(1)}_{\Lambda}]\!]\big)
  \neq 0
  \qquad
  \forall J\in[\Xi^{(1)}_{\Lambda}],
  \]
  and hence a well-defined sign
  \[
  \varepsilon_{\Lambda}
  :=
  \operatorname{sgn}\!\Big(
  \overline{\rho}^{\mathrm{dom}}_{\Lambda}\big([\![\Xi^{(1)}_{\Lambda}]\!]\big)
  \Big)
  \in\{\pm1\};
  \]
\item the normalized orientation functional
  \begin{equation}
  \chi_{\Lambda}
  :=
  \nu_{\Lambda}^{-1}\rho^{\mathrm{dom}}_{\Lambda}
  \label{eq:chi-from-rho-dom}
  \end{equation}
  is well defined, annihilates \(\mathcal A^-_{\Lambda}\), and satisfies
  \[
  \chi_{\Lambda}(J)=\varepsilon_{\Lambda}
  \qquad
  \forall J\in[\Xi^{(1)}_{\Lambda}];
  \]
\item along \(I\), the sign \(\varepsilon_{\Lambda}\) is locally constant once an initial
  orientation is fixed.
\end{enumerate}
\end{theorem}

\begin{proof}
Item (1) is exactly the conclusion of \cref{thm:dominant-line-persistence}.

By assumption, the three generator-wise vanishing statements of
\cref{prop:dominant-quotient-generators} hold for every \(\Lambda\in I\).
Therefore \cref{prop:dominant-quotient-generators,thm:dominant-quotient-descent} give item (2)
and the existence of the descended functional
\(
\overline{\rho}^{\mathrm{dom}}_{\Lambda}
\)
at each cutoff.

Applying \cref{prop:descendant-overlap-equivalent,thm:descendant-overlap} pointwise to the
nonvanishing scalar assumption
\(
\rho^{\mathrm{dom}}_{\Lambda}(J^{\mathrm{desc}}_{\Lambda})\neq 0
\)
now yields
item (3): the entire descendant affine class has one common nonzero dominant-response value, and
its sign \(\varepsilon_{\Lambda}\) is well defined.

Since \(\nu_{\Lambda}>0\) by item (3), the normalized functional \(\chi_{\Lambda}\) of
\eqref{eq:chi-from-rho-dom} is well defined.
Item (2) implies that \(\chi_{\Lambda}\) annihilates \(\mathcal A^-_{\Lambda}\), and item (3)
gives
\[
\chi_{\Lambda}(J)
=
\frac{\rho^{\mathrm{dom}}_{\Lambda}(J)}{
\left|
\overline{\rho}^{\mathrm{dom}}_{\Lambda}\big([\![\Xi^{(1)}_{\Lambda}]\!]\big)
\right|
}
=
\varepsilon_{\Lambda}
\qquad
\forall J\in[\Xi^{(1)}_{\Lambda}],
\]
proving item (4).

Finally, assumption (3) provides a continuous nonvanishing real-valued branchwise descendant scalar,
while items (1)--(3) above identify that scalar with the transported orientation sign of the full
descendant affine class.
Since
\(
\Lambda\mapsto \rho^{\mathrm{dom}}_{\Lambda}(J^{\mathrm{desc}}_{\Lambda})
\)
is continuous and nonzero on the connected set \(I\), its sign is constant, hence in particular
locally constant.
Using item (3), this is exactly the transported orientation sign of the full descendant affine
class, proving item (5).
\end{proof}

\begin{proposition}[Ambiguity-triviality on the compressed welded image]
\label{prop:reduced-welded-ambiguity-trivial}
Assume the hypotheses of \cref{thm:active-response-transport-minimal} at an accepted tick \(k\).
Define the reduced welded image of an odd generator \(F\in E^-_{\Lambda_k}\) by
\[
\mathsf W^{\mathrm{red}}_{\Lambda_k}(F)
:=
\rho^{\mathrm{dom}}_{\Lambda_k}(F)\,P_{-,\Lambda_k}^{\mathrm{LoG}},
\]
where \(P_{-,\Lambda_k}^{\mathrm{LoG}}\) is the canonical compressed odd projector of
\cref{cor:compressed-export-selection}.
Then
\[
\mathsf W^{\mathrm{red}}_{\Lambda_k}(F+A)
=
\mathsf W^{\mathrm{red}}_{\Lambda_k}(F)
\qquad
\forall F\in E^-_{\Lambda_k},\ \forall A\in\mathcal A^-_{\Lambda_k}.
\]
Equivalently, \(\mathsf W^{\mathrm{red}}_{\Lambda_k}\) factors through the ambiguity quotient
\(E^-_{\Lambda_k}/\mathcal A^-_{\Lambda_k}\).
In particular, if \(J_1,J_2\in[\Xi^{(1)}_{\Lambda_k}]\), then
\[
\mathsf W^{\mathrm{red}}_{\Lambda_k}(J_1)
=
\mathsf W^{\mathrm{red}}_{\Lambda_k}(J_2).
\]
\end{proposition}

\begin{proof}
Item (2) of \cref{thm:active-response-transport-minimal} gives
\[
\rho^{\mathrm{dom}}_{\Lambda_k}(A)=0
\qquad
\forall A\in\mathcal A^-_{\Lambda_k}.
\]
Therefore, for every \(F\in E^-_{\Lambda_k}\) and \(A\in\mathcal A^-_{\Lambda_k}\),
\[
\mathsf W^{\mathrm{red}}_{\Lambda_k}(F+A)
=
\rho^{\mathrm{dom}}_{\Lambda_k}(F+A)\,P_{-,\Lambda_k}^{\mathrm{LoG}}
=
\big(\rho^{\mathrm{dom}}_{\Lambda_k}(F)+\rho^{\mathrm{dom}}_{\Lambda_k}(A)\big)
P_{-,\Lambda_k}^{\mathrm{LoG}}
=
\rho^{\mathrm{dom}}_{\Lambda_k}(F)\,P_{-,\Lambda_k}^{\mathrm{LoG}}.
\]
Since \cref{cor:compressed-export-selection} fixes \(P_{-,\Lambda_k}^{\mathrm{LoG}}\) canonically on
the compressed, co-centered, band-locked welded channel, the right-hand side depends only on the
ambiguity class of \(F\).
If \(J_2-J_1\in\mathcal A^-_{\Lambda_k}\), apply the same identity with \(F=J_1\) and
\(A=J_2-J_1\).
\end{proof}

\begin{proposition}[Construction of an admissible calibration for the minimal transported interface]
\label{prop:minimal-transport-calibration}
Fix an accepted tick \(k\). Suppose a weak transported-interface package is given:
\[
\mathcal A^-_{\Lambda_k}\subset E^-_{\Lambda_k},
\qquad
[\Xi^{(1)}_{\Lambda_k}]=J_{0,\Lambda_k}+\mathcal A^-_{\Lambda_k},
\]
together with a continuous nonzero linear functional
\[
\chi_{\Lambda_k}:E^-_{\Lambda_k}\to\mathbb R
\]
and a sign
\(
\varepsilon_{\Lambda_k}\in\{\pm1\}
\)
such that
\[
\mathcal A^-_{\Lambda_k}\subset\ker\chi_{\Lambda_k},
\qquad
\chi_{\Lambda_k}(J)=\varepsilon_{\Lambda_k}
\quad
\forall J\in[\Xi^{(1)}_{\Lambda_k}].
\]
Assume in addition that \(\ker\chi_{\Lambda_k}\) admits a continuous symmetric coercive bilinear form
whose induced norm Hilbertizes it.
Choose any vector \(J^\sharp_{\Lambda_k}\in E^-_{\Lambda_k}\) satisfying
\[
\chi_{\Lambda_k}(J^\sharp_{\Lambda_k})=1,
\]
any scale \(\beta_{\Lambda_k}>0\), and any continuous symmetric coercive bilinear form
\[
\langle\cdot,\cdot\rangle^{\mathrm{aux},\chi}_{\Lambda_k}
:
\ker\chi_{\Lambda_k}\times\ker\chi_{\Lambda_k}\to\mathbb R
\]
whose induced norm Hilbertizes \(\ker\chi_{\Lambda_k}\).
Define
\begin{equation}
\mathsf B_{\Lambda_k}(F,G)
:=
\beta_{\Lambda_k}\,\chi_{\Lambda_k}(F)\chi_{\Lambda_k}(G)
+
\left\langle
F-\chi_{\Lambda_k}(F)J^\sharp_{\Lambda_k},
\,
G-\chi_{\Lambda_k}(G)J^\sharp_{\Lambda_k}
\right\rangle^{\mathrm{aux},\chi}_{\Lambda_k}.
\label{eq:B-minimal-transport}
\end{equation}
Then \(\mathsf B_{\Lambda_k}\) is a continuous symmetric positive definite bilinear form on
\(E^-_{\Lambda_k}\) whose induced norm Hilbertizes \(E^-_{\Lambda_k}\), and
\[
\mathsf B_{\Lambda_k}(J^\sharp_{\Lambda_k},K)=0
\qquad
\forall K\in\ker\chi_{\Lambda_k}.
\]
Hence \(\mathsf B_{\Lambda_k}(J^\sharp_{\Lambda_k},K)=0\) in particular for every
\(K\in\mathcal A^-_{\Lambda_k}\), and
\[
\mathsf B_{\Lambda_k}(J^\sharp_{\Lambda_k},J^\sharp_{\Lambda_k})
=
\beta_{\Lambda_k}.
\]
Moreover, for any representative \(J^\Xi_{\Lambda_k}\in[\Xi^{(1)}_{\Lambda_k}]\), the choice
\[
J^\sharp_{\Lambda_k}:=\varepsilon_{\Lambda_k}J^\Xi_{\Lambda_k}
\]
is admissible.
\end{proposition}

\begin{proof}
Because \(\chi_{\Lambda_k}(J^\sharp_{\Lambda_k})=1\), the linear map
\[
P^\chi_{\Lambda_k}(F)
:=
F-\chi_{\Lambda_k}(F)J^\sharp_{\Lambda_k}
\]
takes values in \(\ker\chi_{\Lambda_k}\):
\[
\chi_{\Lambda_k}\!\big(P^\chi_{\Lambda_k}(F)\big)
=
\chi_{\Lambda_k}(F)-\chi_{\Lambda_k}(F)\chi_{\Lambda_k}(J^\sharp_{\Lambda_k})
=
0.
\]
Thus every \(F\in E^-_{\Lambda_k}\) has the decomposition
\[
F=\chi_{\Lambda_k}(F)J^\sharp_{\Lambda_k}+P^\chi_{\Lambda_k}(F),
\qquad
P^\chi_{\Lambda_k}(F)\in\ker\chi_{\Lambda_k}.
\]
This decomposition is unique: if \(F=aJ^\sharp_{\Lambda_k}+K\) with \(K\in\ker\chi_{\Lambda_k}\),
applying \(\chi_{\Lambda_k}\) gives \(a=\chi_{\Lambda_k}(F)\).

Equation \eqref{eq:B-minimal-transport} is therefore well defined, continuous, and symmetric.
If \(\mathsf B_{\Lambda_k}(F,F)=0\), then
\[
\beta_{\Lambda_k}\chi_{\Lambda_k}(F)^2=0
\]
and
\[
\left\langle
P^\chi_{\Lambda_k}(F),P^\chi_{\Lambda_k}(F)
\right\rangle^{\mathrm{aux},\chi}_{\Lambda_k}
=
0.
\]
Since \(\beta_{\Lambda_k}>0\), we get \(\chi_{\Lambda_k}(F)=0\), hence
\(P^\chi_{\Lambda_k}(F)=F\). Coercivity on \(\ker\chi_{\Lambda_k}\) then implies \(F=0\), proving
positive definiteness.

The map
\[
F\longmapsto
\big(
\chi_{\Lambda_k}(F),P^\chi_{\Lambda_k}(F)
\big)
\]
is a continuous linear isomorphism from \(E^-_{\Lambda_k}\) onto
\(
\mathbb R\oplus\ker\chi_{\Lambda_k}
\),
with inverse
\[
(a,K)\longmapsto aJ^\sharp_{\Lambda_k}+K.
\]
Under this isomorphism, \(\mathsf B_{\Lambda_k}\) becomes the product Hilbert form
\[
(a,K),(b,L)\longmapsto
\beta_{\Lambda_k}ab+\langle K,L\rangle^{\mathrm{aux},\chi}_{\Lambda_k},
\]
so its induced norm Hilbertizes \(E^-_{\Lambda_k}\).

Finally, if \(K\in\ker\chi_{\Lambda_k}\), then \(P^\chi_{\Lambda_k}(J^\sharp_{\Lambda_k})=0\) and
\(\chi_{\Lambda_k}(K)=0\), hence
\[
\mathsf B_{\Lambda_k}(J^\sharp_{\Lambda_k},K)=0.
\]
Setting \(K=J^\sharp_{\Lambda_k}\) gives
\(
\mathsf B_{\Lambda_k}(J^\sharp_{\Lambda_k},J^\sharp_{\Lambda_k})=\beta_{\Lambda_k}
\).
If \(J^\Xi_{\Lambda_k}\in[\Xi^{(1)}_{\Lambda_k}]\), then
\(
\chi_{\Lambda_k}(J^\Xi_{\Lambda_k})=\varepsilon_{\Lambda_k}
\),
so
\(
J^\sharp_{\Lambda_k}:=\varepsilon_{\Lambda_k}J^\Xi_{\Lambda_k}
\)
satisfies
\(
\chi_{\Lambda_k}(J^\sharp_{\Lambda_k})=1
\).
\end{proof}

\begin{remark}[Calibration existence versus stiffness matching]
\label{rem:minimal-transport-calibration-scope}
\Cref{prop:minimal-transport-calibration} is a construction theorem, not an unconditional existence
theorem: it extends a previously chosen Hilbertizing form on \(\ker\chi_{\Lambda_k}\) to an
admissible bilinear form on the whole weak transported interface.
It does \emph{not} identify the scale
\(\beta_{\Lambda_k}\) with the split/leak stiffness \(\widetilde\alpha(\Lambda_k)\).
That physical identification belongs to a later strong or physical layer.
\end{remark}

\begin{theorem}[Rank-one descendant-sector upgrade]
\label{thm:active-response-transport-strong}
Assume \cref{thm:active-response-transport-minimal}, and fix an accepted tick \(k\) on the branch.
Suppose in addition that on the transported
descendant sector the projected response map
\[
F\longmapsto
\Pi^{\mathrm{dom}}_{\Lambda_k}\Pi_{\mathrm{phys},\Lambda_k}\mathcal R_{\Lambda_k}(F),
\qquad
\Pi^{\mathrm{dom}}_{\Lambda_k}
:=
\widehat{\delta C}^{\mathrm{LoG}}_{\Lambda_k}\otimes
\widehat{\delta C}^{\mathrm{LoG}}_{\Lambda_k},
\]
has kernel exactly \(\mathcal A^-_{\Lambda_k}\).
Then the stronger conclusions follow:
\begin{enumerate}[leftmargin=2em]
\item the descendant quotient is rank one,
  \[
  E^-_{\Lambda_k}/\mathcal A^-_{\Lambda_k}\cong \mathbb R;
  \]
\item the orientation kernel is exactly the ambiguity space,
  \[
  \ker\chi_{\Lambda_k}
  =
  \ker\rho^{\mathrm{dom}}_{\Lambda_k}
  =
  \mathcal A^-_{\Lambda_k},
  \]
  so the descendant sector satisfies the rank-one kernel-identification package with quotient
  generated by any vector \(J^{\mathrm{act}}_{\Lambda_k}\) satisfying
  \(\chi_{\Lambda_k}(J^{\mathrm{act}}_{\Lambda_k})=1\);
  \item after choosing any such \(J^{\mathrm{act}}_{\Lambda_k}\), any positive scale
  \(\beta_{\Lambda_k}\), and any auxiliary coercive Hilbertizing form on
  \(\ker\chi_{\Lambda_k}\), \cref{prop:minimal-transport-calibration} yields an admissible
  calibration form.
  Recovering the stiffness-matched formula of \cref{prop:char-calibration-extension} requires
  additional physical identifications, namely
  \(\chi_{\Lambda_k}=\rho_{\Lambda_k}\), \(\ker\chi_{\Lambda_k}=\mathcal T^{\mathrm{op}}_{\Lambda_k}\),
  agreement of the chosen normalized generator with the split/leak active generator, and
  \(\beta_{\Lambda_k}=\widetilde\alpha(\Lambda_k)\).
\end{enumerate}
\end{theorem}

\begin{proof}
By \cref{thm:active-response-transport-minimal}, the dominant line is defined at \(\Lambda_k\), the
scalar \(\nu_{\Lambda_k}\) is positive, and
\[
\chi_{\Lambda_k}
=
\nu_{\Lambda_k}^{-1}\rho^{\mathrm{dom}}_{\Lambda_k}.
\]
Hence
\(
\ker\chi_{\Lambda_k}=\ker\rho^{\mathrm{dom}}_{\Lambda_k}
\).

For every \(F\in E^-_{\Lambda_k}\), the definition of \(\rho^{\mathrm{dom}}_{\Lambda_k}\) gives
\[
\Pi^{\mathrm{dom}}_{\Lambda_k}\Pi_{\mathrm{phys},\Lambda_k}\mathcal R_{\Lambda_k}(F)
=
\rho^{\mathrm{dom}}_{\Lambda_k}(F)\,
\widehat{\delta C}^{\mathrm{LoG}}_{\Lambda_k}.
\]
Therefore the kernel of the projected response map is exactly
\(
\ker\rho^{\mathrm{dom}}_{\Lambda_k}
\).
By the additional hypothesis, that kernel is \(\mathcal A^-_{\Lambda_k}\), proving item (2).

By \cref{thm:active-response-transport-minimal}, there exists
\(
J^{\mathrm{desc}}_{\Lambda_k}\in[\Xi^{(1)}_{\Lambda_k}]
\)
with
\(
\rho^{\mathrm{dom}}_{\Lambda_k}(J^{\mathrm{desc}}_{\Lambda_k})\neq 0
\),
hence the projected response map is nonzero.
Its image is automatically contained in the one-dimensional line
\(
L^{\mathrm{dom}}_{\Lambda_k}
\),
so the first isomorphism theorem gives
\[
E^-_{\Lambda_k}/\mathcal A^-_{\Lambda_k}
\cong
\operatorname{im}\!\big(
\Pi^{\mathrm{dom}}_{\Lambda_k}\Pi_{\mathrm{phys},\Lambda_k}\mathcal R_{\Lambda_k}
\big)
\cong
\mathbb R,
\]
which is item (1).
Any vector \(J^{\mathrm{act}}_{\Lambda_k}\) with \(\chi_{\Lambda_k}(J^{\mathrm{act}}_{\Lambda_k})=1\)
lies outside \(\ker\chi_{\Lambda_k}=\mathcal A^-_{\Lambda_k}\), so its quotient class generates this
one-dimensional quotient, completing item (2).

Item (3) is now an application of \cref{prop:minimal-transport-calibration} with the weak-interface
data supplied by \cref{thm:active-response-transport-minimal} and the chosen normalized vector
\(J^{\mathrm{act}}_{\Lambda_k}\).
The final sentence identifies when this weak calibration construction coincides with the
stiffness-matched formula of \cref{prop:char-calibration-extension}.
\end{proof}

\ifdefined\UOMTransportFragmentLoaded\else
\end{document}
\fi
